%% The Chain, Not the Signature: Cost and Validation of Post-Quantum Content Credentials.
%% Tables and tables/numbers.tex are generated from the released results.

\documentclass[acmsmall,screen,review]{acmart}

\AtBeginDocument{%
  \providecommand\BibTeX{{\normalfont B\kern-0.5em{\scshape i\kern-0.25em b}\kern-0.8em\TeX}}}

\acmJournal{TOPS}
\setcopyright{none}
\settopmatter{printacmref=false}
\renewcommand\footnotetextcopyrightpermission[1]{}

\usepackage{booktabs}
\usepackage{array}
\usepackage{graphicx}
\usepackage{url}
\usepackage{xr}
\externaldocument{supplement}
%% generated by make_tables_v2.py; do not edit
\newcommand{\NarchCrossES}{73}
\newcommand{\NarchCrossHi}{1.4}
\newcommand{\NarchCrossLo}{0.7}
\newcommand{\NarchEffMaxDSAsixtyfive}{0.64}
\newcommand{\NarchEffMaxES}{0.76}
\newcommand{\NarchEffOneDSAsixtyfive}{0.94}
\newcommand{\NarchEffOneES}{0.96}
\newcommand{\NarchErrMaxDSAsixtyfive}{44}
\newcommand{\NarchErrMaxES}{25}
\newcommand{\NarchErrMeanMaxDSAsixtyfive}{36}
\newcommand{\NarchErrMeanMaxES}{24}
\newcommand{\NarchErrMeanOneDSAsixtyfive}{7}
\newcommand{\NarchErrMeanOneES}{4}
\newcommand{\NarchErrOneDSAsixtyfive}{18}
\newcommand{\NarchErrOneES}{5}
\newcommand{\NarchFixMean}{0.58}
\newcommand{\NarchFixMed}{0.50}
\newcommand{\NarchFixPnf}{1.32}
\newcommand{\NarchHeldDSAeightyseven}{0.4}
\newcommand{\NarchHeldDSAfortyfour}{10.3}
\newcommand{\NarchHeldDSAsixtyfive}{2.1}
\newcommand{\NarchIqrRelDSAeightyseven}{0.68}
\newcommand{\NarchIqrRelDSAfortyfour}{0.95}
\newcommand{\NarchIqrRelDSAsixtyfive}{0.92}
\newcommand{\NarchMaxWorkers}{8}
\newcommand{\NarchPiCrossES}{149}
\newcommand{\NarchPiCrossHi}{2.6}
\newcommand{\NarchPiCrossLo}{1.3}
\newcommand{\NarchPiEffMaxDSAsixtyfive}{0.97}
\newcommand{\NarchPiEffMaxDSAsixtyfiveHi}{0.98}
\newcommand{\NarchPiEffMaxDSAsixtyfiveLo}{0.93}
\newcommand{\NarchPiEffMaxES}{0.88}
\newcommand{\NarchPiEffMaxESHi}{0.92}
\newcommand{\NarchPiEffMaxESLo}{0.74}
\newcommand{\NarchPiEffOneDSAsixtyfive}{1.00}
\newcommand{\NarchPiEffOneDSAsixtyfiveHi}{1.03}
\newcommand{\NarchPiEffOneDSAsixtyfiveLo}{0.96}
\newcommand{\NarchPiEffOneES}{0.95}
\newcommand{\NarchPiEffOneESHi}{0.95}
\newcommand{\NarchPiEffOneESLo}{0.95}
\newcommand{\NarchPiErrMaxDSAsixtyfive}{20}
\newcommand{\NarchPiErrMaxDSAsixtyfiveHi}{21}
\newcommand{\NarchPiErrMaxDSAsixtyfiveLo}{19}
\newcommand{\NarchPiErrMaxES}{12}
\newcommand{\NarchPiErrMaxESHi}{26}
\newcommand{\NarchPiErrMaxESLo}{8}
\newcommand{\NarchPiErrMeanAllMaxDSAsixtyfive}{7}
\newcommand{\NarchPiErrMeanAllMaxES}{26}
\newcommand{\NarchPiErrMeanMaxDSAsixtyfive}{3}
\newcommand{\NarchPiErrMeanMaxDSAsixtyfiveHi}{7}
\newcommand{\NarchPiErrMeanMaxDSAsixtyfiveLo}{2}
\newcommand{\NarchPiErrMeanMaxES}{12}
\newcommand{\NarchPiErrMeanMaxESHi}{26}
\newcommand{\NarchPiErrMeanMaxESLo}{8}
\newcommand{\NarchPiErrMeanOneDSAsixtyfive}{3}
\newcommand{\NarchPiErrMeanOneDSAsixtyfiveHi}{4}
\newcommand{\NarchPiErrMeanOneDSAsixtyfiveLo}{0}
\newcommand{\NarchPiErrMeanOneES}{5}
\newcommand{\NarchPiErrMeanOneESHi}{5}
\newcommand{\NarchPiErrMeanOneESLo}{5}
\newcommand{\NarchPiErrOneDSAsixtyfive}{17}
\newcommand{\NarchPiErrOneDSAsixtyfiveHi}{18}
\newcommand{\NarchPiErrOneDSAsixtyfiveLo}{16}
\newcommand{\NarchPiErrOneES}{5}
\newcommand{\NarchPiErrOneESHi}{6}
\newcommand{\NarchPiErrOneESLo}{5}
\newcommand{\NarchPiFixMean}{1.93}
\newcommand{\NarchPiFixMed}{1.56}
\newcommand{\NarchPiFixPnf}{4.73}
\newcommand{\NarchPiHeldDSAeightyseven}{0.5}
\newcommand{\NarchPiHeldDSAfortyfour}{0.3}
\newcommand{\NarchPiHeldDSAsixtyfive}{0.8}
\newcommand{\NarchPiIqrRelDSAeightyseven}{0.90}
\newcommand{\NarchPiIqrRelDSAfortyfour}{1.34}
\newcommand{\NarchPiIqrRelDSAsixtyfive}{0.93}
\newcommand{\NarchPiMaxWorkers}{4}
\newcommand{\NarchPiRuns}{3}
\newcommand{\NarchPiTail}{1.24}
\newcommand{\NarchPiTailES}{1.00}
\newcommand{\NarchPiTailOfDSAeightyseven}{1.17}
\newcommand{\NarchPiTailOfDSAfortyfour}{1.41}
\newcommand{\NarchPiTailOfDSAsixtyfive}{1.24}
\newcommand{\NarchTail}{1.16}
\newcommand{\NarchTailES}{1.01}
\newcommand{\NarchTailOfDSAeightyseven}{1.22}
\newcommand{\NarchTailOfDSAfortyfour}{1.35}
\newcommand{\NarchTailOfDSAsixtyfive}{1.16}
\newcommand{\NauditBoth}{2}
\newcommand{\NauditBothWord}{two}
\newcommand{\NauditCc}{5}
\newcommand{\NauditFlags}{6}
\newcommand{\NauditLib}{36}
\newcommand{\NauditMeanNoSpread}{16}
\newcommand{\NauditN}{40}
\newcommand{\NauditNWord}{forty}
\newcommand{\NauditNoSpread}{20}
\newcommand{\NauditNoStat}{3}
\newcommand{\NauditOlevel}{2}
\newcommand{\NauditOqs}{22}
\newcommand{\NauditOqsAny}{4}
\newcommand{\NauditOqsBuild}{2}
\newcommand{\NauditPreprints}{16}
\newcommand{\NauditSpread}{20}
\newcommand{\NauditVer}{17}
\newcommand{\NavxSign}{2.6}
\newcommand{\NavxVerify}{3.5}
\newcommand{\NbleNjPerByte}{112}
\newcommand{\NbleRsq}{1.0000}
\newcommand{\NbleRxA}{-3}
\newcommand{\NbleRxN}{20}
\newcommand{\NbleRxNj}{115}
\newcommand{\NbleRxNjHi}{115}
\newcommand{\NbleRxNjLo}{114}
\newcommand{\NbleRxResPct}{30}
\newcommand{\NbleRxResUj}{26}
\newcommand{\NbleRxRsq}{0.9999}
\newcommand{\NbleRxXmax}{27,935}
\newcommand{\NbleRxXmin}{113}
\newcommand{\NbleTxA}{-5}
\newcommand{\NbleTxN}{20}
\newcommand{\NbleTxNj}{111}
\newcommand{\NbleTxNjHi}{112}
\newcommand{\NbleTxNjLo}{111}
\newcommand{\NbleTxResPct}{44}
\newcommand{\NbleTxResUj}{14}
\newcommand{\NbleTxRsq}{0.9999}
\newcommand{\NbleTxXmax}{27,935}
\newcommand{\NbleTxXmin}{113}
\newcommand{\NbleWorst}{27.8}
\newcommand{\NcacheLongDSAeightyseven}{2.6}
\newcommand{\NcacheLongDSAfortyfour}{2.6}
\newcommand{\NcacheLongDSAsixtyfive}{2.6}
\newcommand{\NcacheLongES}{3.0}
\newcommand{\NcacheLongEd}{2.9}
\newcommand{\NcacheLongPS}{3.2}
\newcommand{\NcachePerDSAeightyseven}{1.4}
\newcommand{\NcachePerDSAfortyfour}{1.4}
\newcommand{\NcachePerDSAsixtyfive}{1.4}
\newcommand{\NcachePerES}{1.4}
\newcommand{\NcachePerEd}{1.4}
\newcommand{\NcachePerPS}{1.5}
\newcommand{\NcacheRotDSAeightyseven}{2.7}
\newcommand{\NcacheRotDSAfortyfour}{2.3}
\newcommand{\NcacheRotDSAsixtyfive}{2.6}
\newcommand{\NcacheRotES}{2.7}
\newcommand{\NcacheRotEd}{2.7}
\newcommand{\NcacheRotPS}{3.2}
\newcommand{\Ncachepublishers}{50}
\newcommand{\Ncachestream}{2,000}
\newcommand{\NcarrChainDSAeightyseven}{21,419}
\newcommand{\NcarrChainDSAfortyfour}{12,238}
\newcommand{\NcarrChainDSAsixtyfive}{16,185}
\newcommand{\NcarrChainES}{2,658}
\newcommand{\NcarrChainEd}{2,544}
\newcommand{\NcarrChainPS}{3,642}
\newcommand{\NcarrLeafDSAeightyseven}{13,569}
\newcommand{\NcarrLeafDSAfortyfour}{7,875}
\newcommand{\NcarrLeafDSAsixtyfive}{10,293}
\newcommand{\NcarrLeafES}{1,906}
\newcommand{\NcarrLeafEd}{1,849}
\newcommand{\NcarrLeafPS}{2,495}
\newcommand{\NcarrRootDSAeightyseven}{29,007}
\newcommand{\NcarrRootDSAfortyfour}{16,339}
\newcommand{\NcarrRootDSAsixtyfive}{21,815}
\newcommand{\NcarrRootES}{3,155}
\newcommand{\NcarrRootEd}{2,981}
\newcommand{\NcarrRootPS}{4,535}
\newcommand{\NcarrThumbDSAeightyseven}{5,750}
\newcommand{\NcarrThumbDSAfortyfour}{3,543}
\newcommand{\NcarrThumbDSAsixtyfive}{4,432}
\newcommand{\NcarrThumbES}{1,185}
\newcommand{\NcarrThumbEd}{1,185}
\newcommand{\NcarrThumbPS}{1,379}
\newcommand{\NchainBytesDSAeightyseven}{15,697}
\newcommand{\NchainBytesDSAfortyfour}{8,723}
\newcommand{\NchainBytesDSAsixtyfive}{11,781}
\newcommand{\NchainBytesES}{1,501}
\newcommand{\NchainBytesEd}{1,387}
\newcommand{\NchainBytesPS}{2,291}
\newcommand{\NchainShareDSAeightyseven}{77}
\newcommand{\NchainShareDSAfortyfour}{78}
\newcommand{\NchainShareDSAsixtyfive}{78}
\newcommand{\NchainShareES}{95}
\newcommand{\NchainShareEd}{94}
\newcommand{\NchainShareManDSAeightyseven}{73}
\newcommand{\NchainShareManDSAfortyfour}{71}
\newcommand{\NchainShareManDSAsixtyfive}{73}
\newcommand{\NchainShareManES}{56}
\newcommand{\NchainShareManEd}{55}
\newcommand{\NchainShareManPS}{63}
\newcommand{\NchainSharePS}{89}
\newcommand{\NcompositeCose}{15,501}
\newcommand{\NcompressDSAeightyseven}{4}
\newcommand{\NcompressDSAfortyfour}{6}
\newcommand{\NcompressDSAsixtyfive}{5}
\newcommand{\NcompressES}{41}
\newcommand{\NcompressEd}{45}
\newcommand{\NcompressPS}{27}
\newcommand{\NcorpMultiGroup}{8}
\newcommand{\NcorpProdAssets}{317}
\newcommand{\NcorpProdChainMed}{1,865}
\newcommand{\NcorpProdChains}{23}
\newcommand{\NcorpProdEs}{15}
\newcommand{\NcorpProdGens}{28}
\newcommand{\NcorpProdIssuer}{2}
\newcommand{\NcorpProdListed}{9}
\newcommand{\NcorpProdNxfive}{425}
\newcommand{\NcorpProdOcsp}{114}
\newcommand{\NcorpProdOneCert}{3}
\newcommand{\NcorpProdOneCertListed}{0}
\newcommand{\NcorpProdPs}{8}
\newcommand{\NcorpProdPthree}{355}
\newcommand{\NcorpProdRoot}{4}
\newcommand{\NcorpProdSigEs}{270}
\newcommand{\NcorpProdSigPs}{155}
\newcommand{\NcorpProdSigs}{425}
\newcommand{\NcorpProdTs}{352}
\newcommand{\NcorpProdTsaListed}{166}
\newcommand{\NcorpProdTstHi}{6,640}
\newcommand{\NcorpProdTstLo}{1,941}
\newcommand{\NcorpProdTstMed}{3,663}
\newcommand{\NcorpProdTstN}{352}
\newcommand{\NcorpProdTwoCert}{16}
\newcommand{\NcorpProdVtwo}{168}
\newcommand{\NcorpProdVtwoString}{7}
\newcommand{\NcorpTestAssets}{25}
\newcommand{\NcorpTestChainMed}{3,070}
\newcommand{\NcorpTestChains}{4}
\newcommand{\NcorpTestEs}{1}
\newcommand{\NcorpTestGens}{1}
\newcommand{\NcorpTestIssuer}{0}
\newcommand{\NcorpTestListed}{0}
\newcommand{\NcorpTestNxfive}{40}
\newcommand{\NcorpTestOcsp}{0}
\newcommand{\NcorpTestOneCert}{0}
\newcommand{\NcorpTestOneCertListed}{0}
\newcommand{\NcorpTestPs}{3}
\newcommand{\NcorpTestPthree}{0}
\newcommand{\NcorpTestRoot}{1}
\newcommand{\NcorpTestSigs}{41}
\newcommand{\NcorpTestTs}{39}
\newcommand{\NcorpTestTstMed}{5,940}
\newcommand{\NcorpTestTwoCert}{3}
\newcommand{\NcorpTestVtwo}{0}
\newcommand{\NcoseDSAeightyseven}{20,347}
\newcommand{\NcoseDSAfortyfour}{11,166}
\newcommand{\NcoseDSAsixtyfive}{15,113}
\newcommand{\NcoseES}{1,586}
\newcommand{\NcoseEd}{1,472}
\newcommand{\NcosePS}{2,570}
\newcommand{\NcostRootDSAeightyseven}{35}
\newcommand{\NcostRootDSAfortyfour}{34}
\newcommand{\NcostRootDSAsixtyfive}{35}
\newcommand{\NcostRootES}{19}
\newcommand{\NcostRootEd}{17}
\newcommand{\NcostRootPS}{25}
\newcommand{\NcpqAddMaxkB}{550}
\newcommand{\NcpqAddMedkB}{27}
\newcommand{\NcpqAddPNinetykB}{54}
\newcommand{\NcpqFileMedMB}{2.2}
\newcommand{\NcpqFiles}{317}
\newcommand{\NcpqFilesOther}{257}
\newcommand{\NcpqOverOne}{185}
\newcommand{\NcpqOverTen}{80}
\newcommand{\NcpqOverTenOther}{20}
\newcommand{\NcpqOverTenTop}{60}
\newcommand{\NcpqPctMax}{57}
\newcommand{\NcpqPctMed}{1.2}
\newcommand{\NcpqPctPNinety}{16}
\newcommand{\NcpqStoreGrowth}{3.4}
\newcommand{\NdeltaESDSAeightyseven}{18.8}
\newcommand{\NdeltaESDSAfortyfour}{9.6}
\newcommand{\NdeltaESDSAsixtyfive}{13.5}
\newcommand{\NdeltaESES}{0.0}
\newcommand{\NdeltaESEd}{-0.1}
\newcommand{\NdeltaESPS}{1.0}
\newcommand{\NdepthSpread}{11}
\newcommand{\NedgeColdDSAsixtyfive}{0.33}
\newcommand{\NedgeColdEd}{0.20}
\newcommand{\NedgeWarmDiff}{0.04}
\newcommand{\NenergyVthree}{yes}
\newcommand{\NfalconEmbPct}{5.7}
\newcommand{\NfalconRefPct}{1.7}
\newcommand{\NfullGrowthDSAeightyseven}{2.6}
\newcommand{\NfullGrowthDSAfortyfour}{2.5}
\newcommand{\NfullGrowthDSAsixtyfive}{2.6}
\newcommand{\NfullGrowthES}{2.2}
\newcommand{\NfullGrowthEd}{2.2}
\newcommand{\NfullGrowthPS}{2.3}
\newcommand{\NhashHd}{0.365}
\newcommand{\NhashKodak}{0.044}
\newcommand{\NhashPhone}{1.619}
\newcommand{\NhashRaw}{4.915}
\newcommand{\NhashShareHd}{76}
\newcommand{\NhashShareKodak}{28}
\newcommand{\NhashSharePhone}{94}
\newcommand{\NhashShareRaw}{98}
\newcommand{\NhashShareSocial}{56}
\newcommand{\NhashShareThumb}{7}
\newcommand{\NhashSocial}{0.143}
\newcommand{\NhashThumb}{0.009}
\newcommand{\NhwRatioDSAeightyseven}{5.0}
\newcommand{\NhwRatioDSAfortyfour}{1.8}
\newcommand{\NhwSpeedup}{7.6}
\newcommand{\NleafBytesDSAeightyseven}{7,851}
\newcommand{\NleafBytesDSAfortyfour}{4,364}
\newcommand{\NleafBytesDSAsixtyfive}{5,893}
\newcommand{\NleafBytesES}{752}
\newcommand{\NleafBytesEd}{696}
\newcommand{\NleafBytesPS}{1,148}
\newcommand{\Nlibgap}{1,538}
\newcommand{\NlibgapClassHi}{1,631}
\newcommand{\NlibgapClassLo}{1,556}
\newcommand{\Nliboutside}{1,128}
\newcommand{\NlibpadHi}{2,088}
\newcommand{\NlibpadLo}{2,018}
\newcommand{\NmanifestDSAeightyseven}{21,419}
\newcommand{\NmanifestDSAfortyfour}{12,238}
\newcommand{\NmanifestDSAsixtyfive}{16,185}
\newcommand{\NmanifestES}{2,658}
\newcommand{\NmanifestEd}{2,544}
\newcommand{\NmanifestPS}{3,642}
\newcommand{\NmanifestkBDSAeightyseven}{21.4}
\newcommand{\NmanifestkBDSAfortyfour}{12.2}
\newcommand{\NmanifestkBDSAsixtyfive}{16.2}
\newcommand{\NmanifestkBES}{2.7}
\newcommand{\NmanifestkBEd}{2.5}
\newcommand{\NmanifestkBPS}{3.6}
\newcommand{\NmixCheap}{8.4}
\newcommand{\NmixClassical}{1.0}
\newcommand{\NmixFull}{9.5}
\newcommand{\NmixInverted}{10.2}
\newcommand{\NmixLadder}{9.4}
\newcommand{\NmixLeafInt}{7.5}
\newcommand{\NmixLeafMig}{4.2}
\newcommand{\NmixPctFullCheap}{88}
\newcommand{\NmixPctFullClassical}{10}
\newcommand{\NmixPctFullFull}{100}
\newcommand{\NmixPctFullInverted}{107}
\newcommand{\NmixPctFullLadder}{99}
\newcommand{\NmixPctFullLeafInt}{79}
\newcommand{\NmixPctFullLeafMig}{44}
\newcommand{\NmixPctFullPkiFirst}{66}
\newcommand{\NmixPkiFirst}{6.3}
\newcommand{\NmixWholeClassicalEs}{3,651}
\newcommand{\NmixWholeClassicalPq}{17,209}
\newcommand{\NmixWholeFullEs}{17,179}
\newcommand{\NmixWholeFullPq}{30,737}
\newcommand{\NmixWholeLeafIntEs}{13,933}
\newcommand{\NmixWholeLeafIntPq}{27,491}
\newcommand{\NmixWholeLeafMigEs}{8,788}
\newcommand{\NmixWholeLeafMigPq}{22,348}
\newcommand{\NmodelErrMax}{1.3}
\newcommand{\NmodelErrPQ}{1.2}
\newcommand{\NmodelF}{1,093}
\newcommand{\NmodelN}{12}
\newcommand{\NmodelO}{1,195}
\newcommand{\NnaiveLongDSAeightyseven}{255}
\newcommand{\NnaiveLongDSAfortyfour}{253}
\newcommand{\NnaiveLongDSAsixtyfive}{254}
\newcommand{\NnaiveLongES}{258}
\newcommand{\NnaiveLongEd}{268}
\newcommand{\NnaiveLongPS}{280}
\newcommand{\NnaivePerDSAeightyseven}{0}
\newcommand{\NnaivePerDSAfortyfour}{0}
\newcommand{\NnaivePerDSAsixtyfive}{0}
\newcommand{\NnaivePerES}{0}
\newcommand{\NnaivePerEd}{0}
\newcommand{\NnaivePerPS}{0}
\newcommand{\NnaiveRotDSAeightyseven}{10}
\newcommand{\NnaiveRotDSAfortyfour}{55}
\newcommand{\NnaiveRotDSAsixtyfive}{30}
\newcommand{\NnaiveRotES}{40}
\newcommand{\NnaiveRotEd}{30}
\newcommand{\NnaiveRotPS}{24}
\newcommand{\NnaiveTotLongDSAeightyseven}{255}
\newcommand{\NnaiveTotLongDSAfortyfour}{253}
\newcommand{\NnaiveTotLongDSAsixtyfive}{254}
\newcommand{\NnaiveTotLongES}{258}
\newcommand{\NnaiveTotLongEd}{268}
\newcommand{\NnaiveTotLongPS}{280}
\newcommand{\NnaiveTotPerDSAeightyseven}{265}
\newcommand{\NnaiveTotPerDSAfortyfour}{272}
\newcommand{\NnaiveTotPerDSAsixtyfive}{250}
\newcommand{\NnaiveTotPerES}{274}
\newcommand{\NnaiveTotPerEd}{250}
\newcommand{\NnaiveTotPerPS}{247}
\newcommand{\NnaiveTotRotDSAeightyseven}{267}
\newcommand{\NnaiveTotRotDSAfortyfour}{249}
\newcommand{\NnaiveTotRotDSAsixtyfive}{261}
\newcommand{\NnaiveTotRotES}{262}
\newcommand{\NnaiveTotRotEd}{281}
\newcommand{\NnaiveTotRotPS}{282}
\newcommand{\NnrfChainES}{504}
\newcommand{\NnrfChainEShw}{66}
\newcommand{\NnrfSignDSAsixtyfive}{145}
\newcommand{\NnucSignDSAsixtyfive}{155}
\newcommand{\NnucSignStackDSAeightyseven}{105}
\newcommand{\NnucSignStackDSAsixtyfive}{67}
\newcommand{\NnucTwentySignMcycles}{6.8}
\newcommand{\NocspCADSAeightyseven}{4,874}
\newcommand{\NocspCADSAfortyfour}{2,667}
\newcommand{\NocspCADSAsixtyfive}{3,556}
\newcommand{\NocspCAES}{315}
\newcommand{\NocspCAEd}{303}
\newcommand{\NocspCAPS}{505}
\newcommand{\NocspDelDSAeightyseven}{12,756}
\newcommand{\NocspDelDSAfortyfour}{7,062}
\newcommand{\NocspDelDSAsixtyfive}{9,480}
\newcommand{\NocspDelES}{1,097}
\newcommand{\NocspDelEd}{1,030}
\newcommand{\NocspDelPS}{1,684}
\newcommand{\Nouroutside}{1,072}
\newcommand{\NpathICADSAeightyseven}{0.191}
\newcommand{\NpathICADSAfortyfour}{0.086}
\newcommand{\NpathICADSAsixtyfive}{0.126}
\newcommand{\NpathICAES}{0.075}
\newcommand{\NpathICAEd}{0.092}
\newcommand{\NpathICAPS}{0.035}
\newcommand{\NpathRootDSAeightyseven}{0.383}
\newcommand{\NpathRootDSAfortyfour}{0.172}
\newcommand{\NpathRootDSAsixtyfive}{0.250}
\newcommand{\NpathRootES}{0.147}
\newcommand{\NpathRootEd}{0.181}
\newcommand{\NpathRootPS}{0.069}
\newcommand{\NpctimgDSAeightyseven}{20.0}
\newcommand{\NpctimgDSAfortyfour}{11.4}
\newcommand{\NpctimgDSAsixtyfive}{15.1}
\newcommand{\NpctimgES}{2.5}
\newcommand{\NpctimgEd}{2.4}
\newcommand{\NpctimgPS}{3.4}
\newcommand{\NpicoChainDSAeightyseven}{412}
\newcommand{\NpicoChainDSAfortyfour}{161}
\newcommand{\NpicoChainDSAsixtyfive}{256}
\newcommand{\NpicoChainES}{440}
\newcommand{\NpicoChainEd}{355}
\newcommand{\NpicoSignDSAeightyseven}{330}
\newcommand{\NpicoSignDSAsixtyfive}{222}
\newcommand{\NpicoSignStackDSAeightyseven}{120}
\newcommand{\NpicoSignStackDSAsixtyfive}{78}
\newcommand{\Nprofileica}{605}
\newcommand{\Nprofilelisthi}{784}
\newcommand{\Nprofilelistlo}{508}
\newcommand{\Nprofilelistmed}{637}
\newcommand{\NramChainDSAeightyseven}{48}
\newcommand{\NramChainDSAfortyfour}{29}
\newcommand{\NramChainDSAsixtyfive}{37}
\newcommand{\NramChainES}{5}
\newcommand{\NramChainEd}{5}
\newcommand{\NramLeafDSAeightyseven}{32}
\newcommand{\NramLeafDSAfortyfour}{20}
\newcommand{\NramLeafDSAsixtyfive}{25}
\newcommand{\NramLeafES}{3}
\newcommand{\NramLeafEd}{4}
\newcommand{\NramPicoDSAeightyseven}{127}
\newcommand{\NramPicoDSAfortyfour}{55}
\newcommand{\NramPicoDSAsixtyfive}{83}
\newcommand{\NramPicoES}{5}
\newcommand{\NramPicoEd}{5}
\newcommand{\NramRootDSAeightyseven}{63}
\newcommand{\NramRootDSAfortyfour}{37}
\newcommand{\NramRootDSAsixtyfive}{48}
\newcommand{\NramRootES}{6}
\newcommand{\NramRootEd}{6}
\newcommand{\NramThumbDSAeightyseven}{17}
\newcommand{\NramThumbDSAfortyfour}{12}
\newcommand{\NramThumbDSAsixtyfive}{14}
\newcommand{\NramThumbES}{2}
\newcommand{\NramThumbEd}{3}
\newcommand{\NrateDSAsixtyfiveCac}{7,401}
\newcommand{\NrateDSAsixtyfiveOqs}{12,629}
\newcommand{\NrateDSAsixtyfiveUnc}{2,509}
\newcommand{\NrateESCac}{15,429}
\newcommand{\NrateESUnc}{4,432}
\newcommand{\NrateEdCac}{11,231}
\newcommand{\NrateEdUnc}{3,600}
\newcommand{\NratePSCac}{40,446}
\newcommand{\NratePSUnc}{9,481}
\newcommand{\NrecvMChainDSAeightyseven}{2.36}
\newcommand{\NrecvMChainDSAfortyfour}{1.27}
\newcommand{\NrecvMChainDSAsixtyfive}{1.71}
\newcommand{\NrecvMChainES}{0.18}
\newcommand{\NrecvMChainEd}{0.17}
\newcommand{\NrecvMFullDSAeightyseven}{3.19}
\newcommand{\NrecvMFullDSAfortyfour}{1.74}
\newcommand{\NrecvMFullDSAsixtyfive}{2.39}
\newcommand{\NrecvMFullES}{0.24}
\newcommand{\NrecvMFullEd}{0.21}
\newcommand{\NrecvMLeafDSAeightyseven}{1.42}
\newcommand{\NrecvMLeafDSAfortyfour}{0.78}
\newcommand{\NrecvMLeafDSAsixtyfive}{1.05}
\newcommand{\NrecvMLeafES}{0.09}
\newcommand{\NrecvMLeafEd}{0.09}
\newcommand{\NrecvMThumbDSAeightyseven}{0.53}
\newcommand{\NrecvMThumbDSAfortyfour}{0.28}
\newcommand{\NrecvMThumbDSAsixtyfive}{0.38}
\newcommand{\NrecvMThumbES}{0.01}
\newcommand{\NrecvMThumbEd}{0.01}
\newcommand{\NrecvShareVerify}{47}
\newcommand{\NsaveLeafDSAeightyseven}{37}
\newcommand{\NsaveLeafDSAfortyfour}{36}
\newcommand{\NsaveLeafDSAsixtyfive}{36}
\newcommand{\NsaveLeafES}{28}
\newcommand{\NsaveLeafEd}{27}
\newcommand{\NsaveLeafPS}{31}
\newcommand{\NsaveThumbDSAeightyseven}{73}
\newcommand{\NsaveThumbDSAfortyfour}{71}
\newcommand{\NsaveThumbDSAsixtyfive}{73}
\newcommand{\NsaveThumbES}{55}
\newcommand{\NsaveThumbEd}{53}
\newcommand{\NsaveThumbPS}{62}
\newcommand{\NscMCacDSAfortyfour}{11,234}
\newcommand{\NscMCacDSAsixtyfive}{7,401}
\newcommand{\NscMCacES}{15,429}
\newcommand{\NscMCacEd}{11,231}
\newcommand{\NscMCacPS}{40,446}
\newcommand{\NscMCoresCacDSAfortyfour}{1}
\newcommand{\NscMCoresCacDSAsixtyfive}{2}
\newcommand{\NscMCoresCacES}{1}
\newcommand{\NscMCoresCacEd}{1}
\newcommand{\NscMCoresCacPS}{1}
\newcommand{\NscMCoresOqsDSAfortyfour}{1}
\newcommand{\NscMCoresOqsDSAsixtyfive}{1}
\newcommand{\NscMCoresUncDSAfortyfour}{3}
\newcommand{\NscMCoresUncDSAsixtyfive}{4}
\newcommand{\NscMCoresUncES}{3}
\newcommand{\NscMCoresUncEd}{3}
\newcommand{\NscMCoresUncPS}{2}
\newcommand{\NscMMax}{8}
\newcommand{\NscMOqsDSAfortyfour}{19,634}
\newcommand{\NscMOqsDSAsixtyfive}{12,629}
\newcommand{\NscMSpeedCacDSAfortyfour}{4.7}
\newcommand{\NscMSpeedCacDSAsixtyfive}{4.9}
\newcommand{\NscMSpeedCacES}{4.7}
\newcommand{\NscMSpeedCacEd}{4.9}
\newcommand{\NscMSpeedCacPS}{4.2}
\newcommand{\NscMSpeedOqsDSAfortyfour}{4.5}
\newcommand{\NscMSpeedOqsDSAsixtyfive}{4.5}
\newcommand{\NscMSpeedUncDSAfortyfour}{5.8}
\newcommand{\NscMSpeedUncDSAsixtyfive}{5.9}
\newcommand{\NscMSpeedUncES}{5.8}
\newcommand{\NscMSpeedUncEd}{5.9}
\newcommand{\NscMSpeedUncPS}{5.9}
\newcommand{\NscMUncDSAfortyfour}{3,758}
\newcommand{\NscMUncDSAsixtyfive}{2,509}
\newcommand{\NscMUncES}{4,432}
\newcommand{\NscMUncEd}{3,600}
\newcommand{\NscMUncPS}{9,481}
\newcommand{\NscMVsES}{57}
\newcommand{\NscMVsPS}{26}
\newcommand{\NscPCacDSAfortyfour}{3,870}
\newcommand{\NscPCacDSAsixtyfive}{2,603}
\newcommand{\NscPCacES}{4,566}
\newcommand{\NscPCacEd}{6,520}
\newcommand{\NscPCacPS}{7,938}
\newcommand{\NscPCoresCacDSAfortyfour}{3}
\newcommand{\NscPCoresCacDSAsixtyfive}{4}
\newcommand{\NscPCoresCacES}{3}
\newcommand{\NscPCoresCacEd}{2}
\newcommand{\NscPCoresCacPS}{2}
\newcommand{\NscPCoresOqsDSAfortyfour}{2}
\newcommand{\NscPCoresOqsDSAsixtyfive}{2}
\newcommand{\NscPCoresUncDSAfortyfour}{8}
\newcommand{\NscPCoresUncDSAsixtyfive}{12}
\newcommand{\NscPCoresUncES}{8}
\newcommand{\NscPCoresUncEd}{5}
\newcommand{\NscPCoresUncPS}{5}
\newcommand{\NscPMax}{4}
\newcommand{\NscPOqsDSAfortyfour}{8,479}
\newcommand{\NscPOqsDSAsixtyfive}{5,598}
\newcommand{\NscPSpeedCacDSAfortyfour}{3.5}
\newcommand{\NscPSpeedCacDSAsixtyfive}{3.6}
\newcommand{\NscPSpeedCacES}{3.5}
\newcommand{\NscPSpeedCacEd}{3.5}
\newcommand{\NscPSpeedCacPS}{3.4}
\newcommand{\NscPSpeedOqsDSAfortyfour}{3.0}
\newcommand{\NscPSpeedOqsDSAsixtyfive}{3.1}
\newcommand{\NscPSpeedUncDSAfortyfour}{4.0}
\newcommand{\NscPSpeedUncDSAsixtyfive}{4.0}
\newcommand{\NscPSpeedUncES}{4.0}
\newcommand{\NscPSpeedUncEd}{4.0}
\newcommand{\NscPSpeedUncPS}{3.9}
\newcommand{\NscPUncDSAfortyfour}{1,299}
\newcommand{\NscPUncDSAsixtyfive}{886}
\newcommand{\NscPUncES}{1,324}
\newcommand{\NscPUncEd}{2,089}
\newcommand{\NscPUncPS}{2,245}
\newcommand{\NscPVsES}{67}
\newcommand{\NscPVsPS}{39}
\newcommand{\NscXCacDSAfortyfour}{6,515}
\newcommand{\NscXCacDSAsixtyfive}{4,162}
\newcommand{\NscXCacES}{10,138}
\newcommand{\NscXCacEd}{7,369}
\newcommand{\NscXCacPS}{26,648}
\newcommand{\NscXCoresCacDSAfortyfour}{2}
\newcommand{\NscXCoresCacDSAsixtyfive}{3}
\newcommand{\NscXCoresCacES}{1}
\newcommand{\NscXCoresCacEd}{2}
\newcommand{\NscXCoresCacPS}{1}
\newcommand{\NscXCoresOqsDSAfortyfour}{1}
\newcommand{\NscXCoresOqsDSAsixtyfive}{1}
\newcommand{\NscXCoresUncDSAfortyfour}{5}
\newcommand{\NscXCoresUncDSAsixtyfive}{7}
\newcommand{\NscXCoresUncES}{4}
\newcommand{\NscXCoresUncEd}{5}
\newcommand{\NscXCoresUncPS}{2}
\newcommand{\NscXMax}{2}
\newcommand{\NscXOqsDSAfortyfour}{20,006}
\newcommand{\NscXOqsDSAsixtyfive}{13,176}
\newcommand{\NscXSpeedCacDSAfortyfour}{1.9}
\newcommand{\NscXSpeedCacDSAsixtyfive}{1.9}
\newcommand{\NscXSpeedCacES}{1.8}
\newcommand{\NscXSpeedCacEd}{1.9}
\newcommand{\NscXSpeedCacPS}{1.8}
\newcommand{\NscXSpeedOqsDSAfortyfour}{1.6}
\newcommand{\NscXSpeedOqsDSAsixtyfive}{1.8}
\newcommand{\NscXSpeedUncDSAfortyfour}{1.9}
\newcommand{\NscXSpeedUncDSAsixtyfive}{2.0}
\newcommand{\NscXSpeedUncES}{1.9}
\newcommand{\NscXSpeedUncEd}{1.9}
\newcommand{\NscXSpeedUncPS}{2.0}
\newcommand{\NscXUncDSAfortyfour}{2,144}
\newcommand{\NscXUncDSAsixtyfive}{1,454}
\newcommand{\NscXUncES}{2,688}
\newcommand{\NscXUncEd}{2,370}
\newcommand{\NscXUncPS}{6,422}
\newcommand{\NscXVsES}{54}
\newcommand{\NscXVsPS}{23}
\newcommand{\NsendChainDSAeightyseven}{2.28}
\newcommand{\NsendChainDSAfortyfour}{1.25}
\newcommand{\NsendChainDSAsixtyfive}{1.69}
\newcommand{\NsendChainES}{0.17}
\newcommand{\NsendChainEd}{0.16}
\newcommand{\NsendLeafDSAeightyseven}{1.40}
\newcommand{\NsendLeafDSAfortyfour}{0.76}
\newcommand{\NsendLeafDSAsixtyfive}{1.03}
\newcommand{\NsendLeafES}{0.09}
\newcommand{\NsendLeafEd}{0.08}
\newcommand{\NsendMChainDSAeightyseven}{2.27}
\newcommand{\NsendMChainDSAfortyfour}{1.23}
\newcommand{\NsendMChainDSAsixtyfive}{1.67}
\newcommand{\NsendMChainES}{0.17}
\newcommand{\NsendMChainEd}{0.16}
\newcommand{\NsendMFullDSAeightyseven}{3.12}
\newcommand{\NsendMFullDSAfortyfour}{1.68}
\newcommand{\NsendMFullDSAsixtyfive}{2.31}
\newcommand{\NsendMFullES}{0.23}
\newcommand{\NsendMFullEd}{0.20}
\newcommand{\NsendMLeafDSAeightyseven}{1.38}
\newcommand{\NsendMLeafDSAfortyfour}{0.74}
\newcommand{\NsendMLeafDSAsixtyfive}{1.02}
\newcommand{\NsendMLeafES}{0.09}
\newcommand{\NsendMLeafEd}{0.09}
\newcommand{\NsendMThumbDSAeightyseven}{0.52}
\newcommand{\NsendMThumbDSAfortyfour}{0.27}
\newcommand{\NsendMThumbDSAsixtyfive}{0.37}
\newcommand{\NsendMThumbES}{0.01}
\newcommand{\NsendMThumbEd}{0.01}
\newcommand{\NsendThumbDSAeightyseven}{0.52}
\newcommand{\NsendThumbDSAfortyfour}{0.27}
\newcommand{\NsendThumbDSAsixtyfive}{0.37}
\newcommand{\NsendThumbES}{0.01}
\newcommand{\NsendThumbEd}{0.01}
\newcommand{\NsignMsDSAeightyseven}{0.165}
\newcommand{\NsignMsDSAfortyfour}{0.074}
\newcommand{\NsignMsDSAsixtyfive}{0.116}
\newcommand{\NsignMsES}{0.022}
\newcommand{\NsignMsEd}{0.030}
\newcommand{\NsignMsPS}{0.535}
\newcommand{\NslhBigPct}{85}
\newcommand{\NstaleExpNaive}{250}
\newcommand{\NstaleExpTot}{250}
\newcommand{\NstaleRevBoundOne}{15}
\newcommand{\NstaleRevBoundOneH}{21}
\newcommand{\NstaleRevBoundSeven}{119}
\newcommand{\NstaleRevBoundSevenH}{145}
\newcommand{\NstaleRevExpiry}{280}
\newcommand{\NstaleRevNaive}{280}
\newcommand{\NstaleRevTot}{280}
\newcommand{\NstreamCpuMax}{0.69}
\newcommand{\NstreamFrameCpuMax}{11}
\newcommand{\NstreamFrameOneKbps}{3,687.1}
\newcommand{\NstreamFrameOneShare}{30.83}
\newcommand{\NstreamLatMax}{6.9}
\newcommand{\NstreamManFourKbps}{32.9}
\newcommand{\NstreamManFourShare}{0.28}
\newcommand{\NstreamManOneKbps}{131.5}
\newcommand{\NstreamManOneShare}{1.10}
\newcommand{\NstreamManSixKbps}{21.9}
\newcommand{\NstreamManSixShare}{0.18}
\newcommand{\NstreamManTwoKbps}{65.7}
\newcommand{\NstreamManTwoShare}{0.55}
\newcommand{\NstreamPnnMax}{24.4}
\newcommand{\NstreamSamplesMin}{144}
\newcommand{\NstreamSeconds}{288}
\newcommand{\NstreamSixSegs}{48}
\newcommand{\NstreamVsiFourKbps}{7.1}
\newcommand{\NstreamVsiFourShare}{0.06}
\newcommand{\NstreamVsiInit}{21,579}
\newcommand{\NstreamVsiOneKbps}{28.3}
\newcommand{\NstreamVsiOneShare}{0.24}
\newcommand{\NstreamVsiSixKbps}{4.7}
\newcommand{\NstreamVsiSixShare}{0.04}
\newcommand{\NstreamVsiTwoKbps}{14.1}
\newcommand{\NstreamVsiTwoShare}{0.12}
\newcommand{\NsweepShareHi}{86}
\newcommand{\NsweepShareLo}{67}
\newcommand{\NsweepSigFactor}{45}
\newcommand{\NtlAnchors}{32}
\newcommand{\NtlBothLists}{16}
\newcommand{\NtlClaimRoots}{20}
\newcommand{\NtlIcaYearsHi}{5}
\newcommand{\NtlIcaYearsLo}{5}
\newcommand{\NtlIssuing}{12}
\newcommand{\NtlRootYearsHi}{30}
\newcommand{\NtlRootYearsLo}{10}
\newcommand{\NtlRootYearsMed}{25}
\newcommand{\NtokenBytesDSAeightyseven}{20,858}
\newcommand{\NtokenBytesDSAfortyfour}{11,677}
\newcommand{\NtokenBytesDSAsixtyfive}{15,624}
\newcommand{\NtokenBytesES}{2,065}
\newcommand{\NtokenBytesEd}{1,975}
\newcommand{\NtokenBytesPS}{3,045}
\newcommand{\NtsBothLeafB}{20,023}
\newcommand{\NtsBothLeafPct}{37}
\newcommand{\NtsBothLeafkB}{20.0}
\newcommand{\NtsFullCADSAeightyseven}{47,151}
\newcommand{\NtsFullCADSAfortyfour}{26,582}
\newcommand{\NtsFullCADSAsixtyfive}{35,365}
\newcommand{\NtsFullCAES}{5,037}
\newcommand{\NtsFullCAEd}{4,822}
\newcommand{\NtsFullCAPS}{7,192}
\newcommand{\NtsFullDSAeightyseven}{55,033}
\newcommand{\NtsFullDSAfortyfour}{30,977}
\newcommand{\NtsFullDSAsixtyfive}{41,289}
\newcommand{\NtsFullES}{5,819}
\newcommand{\NtsFullEd}{5,549}
\newcommand{\NtsFullPS}{8,371}
\newcommand{\NtsGrowthDSAeightyseven}{2.0}
\newcommand{\NtsGrowthDSAfortyfour}{2.0}
\newcommand{\NtsGrowthDSAsixtyfive}{2.0}
\newcommand{\NtsGrowthES}{1.8}
\newcommand{\NtsGrowthEd}{1.8}
\newcommand{\NtsGrowthPS}{1.8}
\newcommand{\NtsSigDSAeightyseven}{21,419}
\newcommand{\NtsSigDSAfortyfour}{12,238}
\newcommand{\NtsSigDSAsixtyfive}{16,185}
\newcommand{\NtsSigES}{2,657}
\newcommand{\NtsSigEd}{2,544}
\newcommand{\NtsSigPS}{3,642}
\newcommand{\NtsTstDSAeightyseven}{42,277}
\newcommand{\NtsTstDSAfortyfour}{23,915}
\newcommand{\NtsTstDSAsixtyfive}{31,809}
\newcommand{\NtsTstES}{4,722}
\newcommand{\NtsTstEd}{4,519}
\newcommand{\NtsTstPS}{6,687}
\newcommand{\NtsTstkBDSAeightyseven}{42.3}
\newcommand{\NtsTstkBDSAfortyfour}{23.9}
\newcommand{\NtsTstkBDSAsixtyfive}{31.8}
\newcommand{\NtsTstkBES}{4.7}
\newcommand{\NtsTstkBEd}{4.5}
\newcommand{\NtsTstkBPS}{6.7}
\newcommand{\NtsVsESDSAeightyseven}{9.0}
\newcommand{\NtsVsESDSAfortyfour}{5.1}
\newcommand{\NtsVsESDSAsixtyfive}{6.7}
\newcommand{\NtsVsESES}{1.0}
\newcommand{\NtsVsESEd}{1.0}
\newcommand{\NtsVsESPS}{1.4}
\newcommand{\NtsVsPSDSAeightyseven}{6.3}
\newcommand{\NtsVsPSDSAfortyfour}{3.6}
\newcommand{\NtsVsPSDSAsixtyfive}{4.8}
\newcommand{\NtsVsPSES}{0.7}
\newcommand{\NtsVsPSEd}{0.7}
\newcommand{\NtsVsPSPS}{1.0}
\newcommand{\NtsaAnchors}{23}
\newcommand{\NtsaIcaUnderClaim}{7}
\newcommand{\NtsaIcaYearsHi}{15}
\newcommand{\NtsaIcaYearsLo}{5}
\newcommand{\NtsaIssuing}{7}
\newcommand{\NtsaLeafSaveB}{5,898}
\newcommand{\NtsaLeafSavePct}{19}
\newcommand{\NvCompleteDSAeightyseven}{1.09}
\newcommand{\NvCompleteDSAfortyfour}{0.50}
\newcommand{\NvCompleteDSAsixtyfive}{0.73}
\newcommand{\NvCompleteES}{0.46}
\newcommand{\NvCompleteEd}{0.57}
\newcommand{\NvCompleteFactorDSAeightyseven}{2.2}
\newcommand{\NvCompleteFactorDSAfortyfour}{2.3}
\newcommand{\NvCompleteFactorDSAsixtyfive}{2.3}
\newcommand{\NvCompleteFactorES}{2.2}
\newcommand{\NvCompleteFactorEd}{2.2}
\newcommand{\NvCompleteFactorPS}{2.5}
\newcommand{\NvCompletePS}{0.22}
\newcommand{\NvCompletePhoneDSAeightyseven}{2.76}
\newcommand{\NvCompletePhoneDSAfortyfour}{2.16}
\newcommand{\NvCompletePhoneDSAsixtyfive}{2.39}
\newcommand{\NvCompletePhoneES}{2.12}
\newcommand{\NvCompletePhoneEd}{2.23}
\newcommand{\NvCompletePhonePS}{1.89}
\newcommand{\NvPathDSAeightyseven}{0.380}
\newcommand{\NvPathDSAfortyfour}{0.170}
\newcommand{\NvPathDSAsixtyfive}{0.251}
\newcommand{\NvPathES}{0.148}
\newcommand{\NvPathEd}{0.181}
\newcommand{\NvPathPS}{0.069}
\newcommand{\NvPathShareDSAeightyseven}{35}
\newcommand{\NvPathShareDSAfortyfour}{34}
\newcommand{\NvPathShareDSAsixtyfive}{34}
\newcommand{\NvPathShareES}{33}
\newcommand{\NvPathShareEd}{32}
\newcommand{\NvPathSharePS}{32}
\newcommand{\NvPathShareSigDSAeightyseven}{77}
\newcommand{\NvPathShareSigDSAfortyfour}{79}
\newcommand{\NvPathShareSigDSAsixtyfive}{78}
\newcommand{\NvPathShareSigES}{72}
\newcommand{\NvPathShareSigEd}{69}
\newcommand{\NvPathShareSigPS}{76}
\newcommand{\NvSigDSAeightyseven}{0.110}
\newcommand{\NvSigDSAfortyfour}{0.043}
\newcommand{\NvSigDSAsixtyfive}{0.068}
\newcommand{\NvSigES}{0.056}
\newcommand{\NvSigEd}{0.081}
\newcommand{\NvSigPS}{0.021}
\newcommand{\NvTotalDSAeightyseven}{0.493}
\newcommand{\NvTotalDSAfortyfour}{0.215}
\newcommand{\NvTotalDSAsixtyfive}{0.322}
\newcommand{\NvTotalES}{0.206}
\newcommand{\NvTotalEd}{0.264}
\newcommand{\NvTotalPS}{0.091}
\newcommand{\NvalAllFive}{37}
\newcommand{\NvalBigCertMB}{44}
\newcommand{\NvalBigCertMs}{16}
\newcommand{\NvalCommit}{e2ad88e}
\newcommand{\NvalPqTsaMs}{2.4}
\newcommand{\NvalRows}{396}
\newcommand{\NvalShown}{24}
\newcommand{\NvalVerifyHi}{3.5}
\newcommand{\NvalVerifyLo}{2.0}
\newcommand{\Nvalcases}{40}
\newcommand{\NverENrfDSAeightyseven}{5.00}
\newcommand{\NverENrfDSAfortyfour}{1.81}
\newcommand{\NverENrfDSAsixtyfive}{3.01}
\newcommand{\NverENrfES}{5.82}
\newcommand{\NverENrfEd}{1.35}
\newcommand{\NverENucDSAeightyseven}{19.82}
\newcommand{\NverENucDSAfortyfour}{7.21}
\newcommand{\NverENucDSAsixtyfive}{11.84}
\newcommand{\NverENucES}{24.95}
\newcommand{\NverENucEd}{6.48}
\newcommand{\NverEPicoDSAeightyseven}{41.65}
\newcommand{\NverEPicoDSAfortyfour}{16.27}
\newcommand{\NverEPicoDSAsixtyfive}{25.87}
\newcommand{\NverEPicoES}{44.90}
\newcommand{\NverEPicoEd}{36.41}
\newcommand{\NvsESDSAeightyseven}{8.1}
\newcommand{\NvsESDSAfortyfour}{4.6}
\newcommand{\NvsESDSAsixtyfive}{6.1}
\newcommand{\NvsESES}{1.0}
\newcommand{\NvsESEd}{1.0}
\newcommand{\NvsESPS}{1.4}
\newcommand{\NvsPSDSAeightyseven}{5.9}
\newcommand{\NvsPSDSAfortyfour}{3.4}
\newcommand{\NvsPSDSAsixtyfive}{4.4}
\newcommand{\NvsPSES}{0.7}
\newcommand{\NvsPSEd}{0.7}
\newcommand{\NvsPSPS}{1.0}
\newcommand{\NxOsslVsOqs}{5.8}
\newcommand{\NxSpread}{6.0}


\begin{document}

\title{The Chain, Not the Signature: Cost and Validation of Post-Quantum
Content Credentials}

\author{Mua'ad Alfarajat}
\orcid{0000-0003-2489-4857}
\affiliation{%
  \institution{The University of Jordan}
  \department{Department of Computer Information Systems,
              School of Information Technology and Systems}
  \city{Aqaba}
  \postcode{77110}
  \country{Jordan}}
\email{m.abufaraj@ju.edu.jo}

\renewcommand{\shortauthors}{Alfarajat}

\begin{abstract}
Cameras, phones and editing tools now sign media with C2PA content
credentials, and every signature a credential rests on becomes forgeable once
a cryptographically relevant quantum computer exists. C2PA admitted ML-DSA in
October 2026. We ask what a conformant post-quantum credential costs, which
choices the specification permits recover it without
weakening what a verifier establishes, and what a migrated credential still
exposes. Certificates, not signatures, carry the cost: before timestamping, an
ML-DSA-65 manifest is \NvsESDSAsixtyfive{} times a matched ES256 one, and
certificates are \NchainShareManDSAsixtyfive\% of it. Substituting measured
ML-DSA-65 components into \NcpqFiles{} production-signed files gives an
estimated median growth of \NcpqAddMedkB\,kB. Because both C2PA trust lists
anchor issuing CAs, a signer or time-stamping authority under a listed issuer
can carry its certificate alone, which takes \NtsBothLeafPct\% off a
timestamped manifest; the reference implementation accepts this form, pending
confirmation by the specification's editors. Partial migration exposes
distinct properties: a classical issuer, or a classical anchor still listed for
the same subject, permits signer impersonation, while a classical
time-stamping authority leaves the signing time forgeable. A conformance
matrix through the reference implementation, patched for ML-DSA, finds two
implementation defects, present with classical algorithms too: it accepts
time-stamping and OCSP-signing certificates as trusted claim signers under a
claim-signing-only policy, and it invalidates a credential over a timestamp it
should ignore. It also ignores revocation stapled in a CA-signed OCSP response,
a gap between the C2PA profile and RFC~6960. We give guidance for CAs,
generators, validators and the specification editors.
\end{abstract}

\begin{CCSXML}
<ccs2012>
<concept>
<concept_id>10002978.10002979.10002981.10011602</concept_id>
<concept_desc>Security and privacy~Digital signatures</concept_desc>
<concept_significance>500</concept_significance>
</concept>
<concept>
<concept_id>10002978.10002979.10002980</concept_id>
<concept_desc>Security and privacy~Key management</concept_desc>
<concept_significance>300</concept_significance>
</concept>
</ccs2012>
\end{CCSXML}
\ccsdesc[500]{Security and privacy~Digital signatures}
\ccsdesc[300]{Security and privacy~Key management}

\keywords{post-quantum migration, content provenance, C2PA, ML-DSA, X.509,
certificate validation, media authentication}

\maketitle


%% ---------------------------------------------------------------------------
\section{Introduction}
\label{sec:intro}

Content provenance has left the proposal stage. Cameras sign at
capture~\cite{cai2023leica,sony2024c2pa}, phones attach content credentials of
their own~\cite{samsung2025s25,google2025pixel}, and editing tools record what
they changed in a signed C2PA manifest that travels inside the file. Each
credential rests on a COSE signature over a claim about the content and on an
X.509 certificate path from the signer to an anchor on the C2PA trust list,
and every signature in it, on the claim, on the certificates and on the
timestamp, is RSA or ECDSA. A cryptographically relevant quantum computer
breaks all of them~\cite{shor1997}, and the PKI being built for credentials
today will still be in service when one arrives: the roots on the C2PA trust list are issued for
a median of \NtlRootYearsMed{} years, most of them into the 2040s and 2050s.

The replacements exist. NIST standardised ML-DSA and SLH-DSA in
2024~\cite{fips204,fips205}; RFC~9964 assigned the COSE identifiers that let an
ML-DSA signature appear in a manifest~\cite{rfc9964} and RFC~9881 those for
ML-DSA certificates~\cite{rfc9881}; and in October 2026 version~2.5 of the
C2PA specification admitted ML-DSA-44, 65 and 87~\cite{c2pa25}. The reference
implementation had not followed by the time of writing~\cite{c2pars}. To our
knowledge nobody has measured what the change does to a credential, and the published
discussion of post-quantum image authenticity proposes new
schemes~\cite{iavich2025pqveritas,alkhatib2025deepfake} rather than pricing the
one that ships.

\paragraph{Research question}
This paper asks what admitting ML-DSA costs a conformant C2PA content
credential, where that cost sits, which choices the specification already
allows reduce it without weakening what a verifier can establish, and what a
migrated credential still exposes. We break the question into three parts.
\begin{itemize}
\item[\textbf{RQ1}] How large and how slow is a conformant post-quantum
credential, timestamped and with revocation evidence, and what does it add to
the files that signers actually publish?
\item[\textbf{RQ2}] Which carriage, trust-list and verifier choices recover
that cost, on servers and on constrained devices, and what does each give up?
\item[\textbf{RQ3}] What does a partly or fully migrated credential still
expose: which partial migrations leave which properties forgeable, and does the
reference implementation enforce what a migration relies on, that time-stamping keys cannot attribute content, that
timestamps are judged as the specification says, and that revoked signers stay
revoked?
\end{itemize}

Size is a security question here, not only a performance one. A credential
that outgrows the space a generator reserves forces a second timestamp request
on every capture; a verifier that cannot hold the chain is pushed towards
thumbprints, which an offline verifier cannot resolve; and a migration that
moves the leaf but not its issuer leaves the content forgeable, while one that
moves the signer but not its time-stamping authority leaves the signing time
forgeable, each a distinct property (Table~\ref{tab:properties}). Which of the
cheaper options stay sound is what the three questions together decide. We measure
application processors, a single-board computer and microcontrollers; the
secure elements in which phones sign are outside what we could measure, and
Section~\ref{sec:constrained} discusses them.

\paragraph{What is known, and what is new}
That certificate chains, not handshake signatures, dominate post-quantum
authentication is established for TLS~\cite{sikeridis2020ndss,paquin2020,%
sosnowski2023,westerbaan2021}, and the TLS remedies avoid sending or
signing what the verifier could already hold or
check~\cite{kampanakis2022scas,certabridge,merkletreecerts,schwabe2020kemtls}.
We confirm the structural result for manifests and do not claim it. What is
new is what changes when the signed object is a file rather than a connection.
A manifest pays for its chain once per object and keeps it for the life of the
file. Nothing is negotiated, so the carriage form is a policy every verifier
must accept. Some verifiers are offline for good. And C2PA runs its own trust
lists, whose anchors include issuing CAs as well as roots, so a signer under a
listed issuer can omit its intermediate with no new mechanism, under a reading
of the specification that the reference implementation already accepts and
whose confirmation we have requested. These differences decide which remedies
carry over. We measure them, and what the reference implementation accepts.
The closest analysis of C2PA itself, Golaszewski et al.'s study of
version~2.2~\cite{golaszewski2026eprint}, finds timestamp disagreement, weak
revocation policy and inconsistent validators, but considers neither
post-quantum signatures nor what a credential costs;
Section~\ref{sec:related} sets out what each covers.

\paragraph{Contributions}
\begin{itemize}
\item \emph{The cost, on real files (RQ1).} A measurement of ML-DSA, SLH-DSA,
Falcon and composite credentials against the ES256 and PS256 credentials
signers ship, in bytes, signing and validation time, with and without a
timestamp and stapled revocation evidence. A closed-form model reproduces the
measured manifests of \NmodelN{} schemes to within \NmodelErrMax\%, so the
sizes extend to other schemes by substitution, and substituting our measured
ML-DSA-65 components into \NcpqFiles{} production-signed files from
\NcorpProdGens{} signing applications estimates what the migration adds to
real files (Section~\ref{sec:cost}).
\item \emph{What recovers the cost (RQ2).} Leaf-only carriage for signers and,
to our knowledge for the first time, for time-stamping authorities, checked
against both live C2PA trust lists, the production chains and the reference
implementation, with what it gives up; how a cache of validated signers must
be keyed to stay sound; and what the carriage form costs constrained
verifiers in memory and radio energy, and what each live-video method costs a
constrained signer (Sections~\ref{sec:carrying} and~\ref{sec:edge}).
\item \emph{What a migrated credential exposes, and what the validator
enforces (RQ3).} A migration-order analysis showing which property each
partial migration leaves forgeable (Table~\ref{tab:properties}), and a
conformance matrix of \Nvalcases{} cases through the reference implementation,
patched to sign and validate ML-DSA, including a fully post-quantum credential
with a post-quantum timestamp. The matrix finds two implementation defects,
time-stamping and OCSP-signing certificates accepted as claim signers and a
timestamp the specification says to ignore invalidating the whole manifest,
and one gap between the C2PA profile and RFC~6960: revocation stapled in a
CA-signed OCSP response has no effect (Section~\ref{sec:findings}).
\end{itemize}
Section~\ref{sec:discussion} turns the results into recommendations for CAs,
claim generators, validators and the specification. Every script, raw
measurement, firmware image and patch is released, as is every generated
certificate except the withheld vectors of Section~\ref{sec:validation}, and
every number in the text is generated from the released measurements.

%% ---------------------------------------------------------------------------
\section{Background}
\label{sec:background}

\subsection{Content credentials}

A C2PA manifest is a record embedded in a media file that states who produced
it, with what device or software, and what was done to it since~\cite{c2pa25}.
Its assertions include a hard binding, a cryptographic hash over the content
that ties the claim to specific bytes. A signature over the claim makes the
record verifiable, and any change to the content breaks the binding.

The signature is a COSE \texttt{Sign1} structure~\cite{rfc9052} with a
detached payload. Its protected header carries the algorithm identifier and the
signer's certificates in the \texttt{x5chain} header of
RFC~9360~\cite{rfc9360}. C2PA narrows RFC~9360 in two ways that matter here.
The claim generator ``shall include the signer's certificate and all intermediate certificate authorities'' in \texttt{x5chain}, and ``the trust anchor's certificate (also called the root certificate) should not be included''; from version~2.5 it ``shall use only the integer 33'' as the label, in the protected bucket, while validators still accept the older string label
\texttt{x5chain} in either bucket~\cite{c2pa25}. A validator trusts a signer
whose certificate path ends at an anchor on the C2PA trust list, a list of
certificates that may be roots or issuing CAs~\cite{c2patrustlist}; a separate
list does the same for time-stamping authorities. In a JPEG the manifest
travels in APP11 segments, so the cost of a credential is the growth in file
size.

Version~2.4 of the specification, on which the measurements were planned,
permitted seven classical algorithms. Version~2.5, published in October 2026,
added ML-DSA-44, 65 and 87 under the COSE identifiers $-48$, $-49$ and
$-50$~\cite{c2pa25}. It adds nothing for SLH-DSA or Falcon.

\subsection{Post-quantum signatures}

ML-DSA~\cite{fips204} is the module-lattice signature scheme NIST
standardised in 2024, at three parameter sets. ML-DSA-44 targets NIST security
category~2 rather than category~1, so it is not the drop-in counterpart of a
128-bit classical parameter set that its name suggests. SLH-DSA~\cite{fips205}
is a hash-based alternative with conservative assumptions and much larger
signatures. FN-DSA, the standard NIST plans for Falcon as FIPS~206, had not
appeared even in draft by 23~October 2026; in September 2026 NIST asked for
comment on a revised plan that specifies fixed-point arithmetic for key
generation and signing~\cite{nist_fips206_status,nist_fips206_forum,%
nist_fips206_plan}. Composite signatures pair ML-DSA with a classical scheme: a verifier accepts only if both component signatures verify, so a forger must break both, and the signature stays secure as long as either algorithm does; the IETF specification is close to
publication~\cite{ounsworth2026composite}, and no COSE identifiers have been
assigned for them.

Signature and public-key sizes are what matter here, because both enter a
manifest: the signature once, and the public key and the issuer's signature
once per certificate. ML-DSA-44, 65 and 87 produce signatures of 2,420, 3,309
and 4,627 bytes and public keys of 1,312, 1,952 and 2,592, against 64 and 32
for Ed25519 and 256 and about 270 for RSA-2048. Certificates for ML-DSA and
SLH-DSA are specified by RFCs~9881 and~9909~\cite{rfc9881,rfc9909}, and CMS
signatures, which a timestamp token is, by RFCs~9882 and~9814~\cite{rfc9882,%
rfc9814}. The IANA COSE registry has no entry for FN-DSA or SLH-DSA; the
working-group draft for FN-DSA requested $-54$ and $-55$, but it expired on
16~September 2026 with both values unassigned~\cite{iana_cose,coseFalcon}.

\subsection{Carrying certificates, or not}

RFC~9360 gives a COSE structure two ways to identify its
signer~\cite{rfc9360}: \texttt{x5chain} carries certificates and \texttt{x5t}
carries a thumbprint, leaving the verifier to find the certificate elsewhere.
Within \texttt{x5chain} there is a further choice of how many certificates to
carry. Section~\ref{sec:carrying} measures four forms: leaf, intermediate and
root; leaf and intermediate, the form C2PA prescribes; the leaf alone; and the
thumbprint.

%% ---------------------------------------------------------------------------
\section{Method}
\label{sec:method}

\subsection{Images, platforms and libraries}

The image corpus is the 24-image Kodak set~\cite{kodak} at JPEG quality 90,
77 to 188\,kB. Kodak images are far smaller than a modern camera produces, and
Section~\ref{sec:denominator} reports cost against file size rather than
against them alone.

The platforms in Table~\ref{tab:platforms} span where content credentials are
signed and verified. A creator's own machine signs at capture or export
(Apple M1 Pro). A cloud service signs and verifies on behalf of many users
inside a virtual machine (a two-core Intel Xeon VM). A camera or edge encoder
signs under a power budget, for which a Raspberry Pi~5 stands in: four
Cortex-A76 cores at a fixed 2.4\,GHz under the performance governor, Raspberry
Pi OS on Debian~13, Linux~6.18, with the work pinned to one core and the
telemetry on the others. Its power-management IC reports voltage and current
per supply rail, which we sum once a second, so its figures are rail power
rather than wall power. Battery-powered verifiers run on microcontrollers; the
three of Section~\ref{sec:mcu} are described there. A figure from one platform
does not transfer to another.

\begin{table}[t]
\caption{Measurement platforms. \emph{Crypto features} lists the hardware instructions the CPU advertises. \emph{OpenSSL} is the optimisation level recorded in the OpenSSL that \texttt{cryptography} links, read from the compiler line OpenSSL embeds in itself; Section~\ref{sec:supp-impl} of the supplement explains why it is reported.}
\label{tab:platforms}
\small
\begin{tabular}{p{3.4cm}lp{2.6cm}p{3.2cm}l}
\toprule
CPU & Arch. & Crypto features & Software & OpenSSL\\
\midrule
Apple M1 Pro & arm64 & aes, sha2, sha3, sha512 & Python 3.12.14, cryptography 50.0.1 & \texttt{-O3}\\
Intel(R) Xeon(R) Processor @ 2.10GHz & x86\_64 & aes, vaes, pclmulqdq, sha\_ni, avx2, avx512f & Python 3.13.17, cryptography 50.0.1 & \texttt{-O3}\\
Raspberry Pi 5 (Cortex-A76) & arm64 & aes, pmull, sha1, sha2 & Python 3.13.5, cryptography 50.0.1 & \texttt{-O3}\\
\bottomrule
\end{tabular}
\end{table}


Signatures run through two libraries: the Python \texttt{cryptography}
package~50.0.1 over OpenSSL~4.0.2, and liboqs~0.16.0. On the M1 Pro the
package is built from source against an OpenSSL compiled at \texttt{-O3},
because the macOS wheel's OpenSSL was compiled without optimisation
(Section~\ref{sec:backends}); on Linux the wheel's OpenSSL records
\texttt{-O3}. Every result file records the compiler line of the OpenSSL it
was measured through.

\subsection{Decisions that shape every number}
\label{sec:method-decisions}

\emph{The chain omits the root.} Every manifest carries the signer's certificate and its intermediate in \texttt{x5chain}, which C2PA~2.5 requires, and leaves out the root, as it says generators should; the root is the validator's trust anchor. Hierarchies are three
certificates deep, the depth deployed signers use
(Section~\ref{sec:prodchains}). Where a table carries the root as well, it
says so.

\emph{The classical baseline is what deployments ship.} Production signers
use ES256 about twice as often as RSA-2048 PS256
(Section~\ref{sec:prodchains}), so both appear in every table beside Ed25519,
and growth is reported against both.

\emph{Every certificate is real and follows a production profile.} Names
carry organisation and country; the leaf carries the
\texttt{c2pa-kp-claimSigning} extended key usage that C2PA~2.5 names; every
certificate below the root carries authority information access, a CRL
distribution point and a certificate policy, as the issuing CAs on the C2PA
trust list do. Our issuing-CA certificate carries \Nprofileica{} bytes beyond
its key and signature, against \Nprofilelistlo{} to \Nprofilelisthi{}
(median \Nprofilelistmed{}) for the issuing CAs on the trust list. ML-DSA and
classical certificates come from \texttt{cryptography}. No library we could
use issues Falcon, SLH-DSA or composite certificates, so we encode the same
profile and sign it with liboqs or with both component keys; the encoder
reproduces \texttt{cryptography}'s ML-DSA certificates byte for byte, and
every chain verifies. FN-DSA and the composites have no final object
identifiers, so those certificates carry placeholders of the same length.

\emph{Timestamp tokens and OCSP responses are real.} No library available to
us issues an RFC~3161 token under a post-quantum key, so we build each one: a
CMS SignedData holding the TSTInfo and the signed attributes the standard
requires, signed by a time-stamping authority whose certificate and
intermediate travel in the token. OCSP responses are real BasicOCSPResponse
objects in both forms RFC~6960 allows, signed by the issuing CA itself or by a
delegated responder whose certificate travels in the response.

\emph{Thumbnails are disabled and reservation is reported separately.} A C2PA
manifest embeds a generated thumbnail by default; on the first Kodak image it
ran to 173\,kB, larger than the photograph. The reference implementation also pads the signature to a reserved size, which its default classical signer set at about 12\,kB in the version we first measured (0.90.19, Table~\ref{tab:baseline} in the supplement). Both would swamp every signature effect,
so sizes here are payload sizes. The reservation still matters: when it is too
small, the specification has the signer start again with more padding and a
fresh timestamp~\cite{c2pacc24}, so a signer that keeps a classical
reservation and moves to ML-DSA-65, whose signed structure is
\NtsSigDSAsixtyfive{} bytes, pays a second round trip to a time-stamping
authority on every capture.

\emph{Falcon signatures are padded.} Falcon-512 produced signatures between
648 and 661 bytes over 200 trials; a reservation cannot be variable, so the
tables use the padded 666 bytes (1,280 for Falcon-1024).

\subsection{Timing and statistics}
\label{sec:timing}

The size and timing experiments ran ten times on the M1 Pro and the Xeon VM (nine for the Xeon's cache simulation), and verification, hashing and carriage five times on the Pi~5; the scaling experiment ran three times on each machine. Each run is a fresh set of processes with fresh keys, its own
random order of experiments and, within each experiment, of algorithms and
configurations. Within a run, a signature timing is taken over 1,000
operations after 50 discarded, and path validation and hashing over 200. ML-DSA
signing in both libraries is hedged, so repeating one message still draws a
fresh rejection-sampling path each time. We report the median of the per-run
medians with a 95\% bootstrap interval over runs; within-run percentiles are
in the released data, and the 99th percentile is quoted where a real-time
claim rests on it. The replication unit is therefore the run, not the repetition. The stream and edge experiments on the Pi ran once, with their repetitions inside the run, and the archive batch three times. The cache-staleness and composite-rejection tests count outcomes rather than time, and ran on the analysis host. ECDSA signatures inside certificates vary in DER length from key to key, so ES256 sizes differ by a byte or two between experiments that each generate their own keys. Signature timings are taken over the COSE \texttt{Sig\_structure}, the exact bytes a signature covers, so libraries are compared on identical input; the algorithm sweep alone signs the claim, which leaves out the certificates in the protected header and makes its ML-DSA times 16 to 18 per cent lower.

\subsection{Through the reference implementation}
\label{sec:patch}

Most structures in this paper are built directly to RFC~9964, which leaves a
question: does a real claim generator, doing all the work a claim generator
does, produce and accept what we say it does? Settling it meant patching the
reference implementation, \texttt{c2pa-rs}, which signs only classical
algorithms, and the patch is itself a finding about how ready the tooling
is.

\texttt{coset}, the COSE library underneath, already carries the RFC~9964
identifiers. Everything missing sat in the layer above, behind four gates
that each had to be opened: the signing-algorithm enumeration, which needed
three variants and an ML-DSA signer and validator of about 140 lines against
the RustCrypto \texttt{ml-dsa} crate; the COSE header mapping in each
direction, whose catch-all arms reject without a compiler error; the
verifier's check of the key's object identifier; and the certificate
profile's own allow-list of certificate signature algorithms. That patch, made
against the commit of 24~September, applies unchanged to the current
\texttt{main} (commit \texttt{\NvalCommit}). On that commit it is no longer
enough. Path validation now maps each certificate's algorithm to a validator
of its own, with no entry for ML-DSA, so every ML-DSA chain validates as
\emph{Valid} and never as \emph{Trusted}; the OCSP check has the same gap. A
second patch of 82 lines adds the mappings.

\begin{table*}[t]
\caption{Validation through the reference implementation (\texttt{c2pa-rs} \texttt{e2ad88e} with our ML-DSA patches, p2): a selection of the 40 cases; Table~\ref{tab:valfull} in the supplement lists all. T, V, I: Trusted, Valid, Invalid; E: reader error; SR: signing refused (\emph{direct}: our harness's own structure with that certificate). \emph{Expected} follows C2PA~2.5. \emph{TST}: the timestamp's status (t trusted, m mismatch, u untrusted); an ignored timestamp leaves the signer judged now. The TSA is ES256 except in the indented rows. $^\dagger$Bold: a defect, also present unpatched. The CA-signed \emph{revoked} OCSP row is where C2PA~2.5 and RFC~6960 disagree (Section~\ref{sec:validation}).}
\label{tab:validation}
\footnotesize
\setlength{\tabcolsep}{3pt}
\begin{tabular}{llccc}
\toprule
 & & ES256, & ML-DSA & \\
Case & Expected & Ed25519 & 44/65/87 & TST\\
\midrule
Leaf and intermediate & T & T & T & --\\
\quad timestamped & T & T & T & t\\
Leaf, intermediate, root & T & T & T & t\\
Leaf only, issuer is anchor & T & T & T & t\\
Leaf only, intermediate missing & not T & V & V & t\\
Unknown anchor & not T & V & V & t\\
Thumbprint (\texttt{x5t}) & not T & E & E & --\\
Leaf EKU serverAuth & refused & SR (direct I) & SR (direct I) & -- (direct --)\\
Leaf EKU timeStamping only$^\dagger$ & not T & \textbf{T} & \textbf{T} & t\\
COSE algorithm $\neq$ key (ML-DSA-44 only) & I & I & I & --\\
Image changed after signing & I & I & I & t\\
Signature bit flipped & I & I & I & m\\
Expired leaf, timestamp inside validity & T & T & T & t\\
Expired leaf, no timestamp & not T & I & I & --\\
Expired leaf, timestamp over other data & not T & I & I & m\\
Expired leaf, TSA not trusted & not T & I & I & u\\
Timestamp over other data & T & T & T & m\\
TSA not trusted & T & T & T & u\\
TSA certificate lacks timeStamping$^\dagger$ & T & \textbf{I} & \textbf{I} & u\\
ML-DSA-65 TSA & T & T & T & t\\
\quad timestamp over other data & T & T & T & m\\
\quad expired leaf, timestamp inside validity & T & T & T & t\\
OCSP, CA-signed, revoked & T & T & T & t\\
OCSP, delegated responder, revoked & not T & I & I & t\\
\bottomrule
\end{tabular}
\end{table*}


We ran \Nvalcases{} cases through the patched library, \NvalAllFive{} of them for each of ES256, Ed25519 and the three ML-DSA parameter sets and the rest for the algorithms they concern, \NvalRows{} validations in all: the four carriage forms, the negative cases a validator must reject, timestamps from an ES256 and from an ML-DSA-65 time-stamping authority, both OCSP forms, expiry with good and with rejected timestamps, and malformed and oversized structures. Table~\ref{tab:validation} shows \NvalShown{} of them and the supplement lists every case (Table~\ref{tab:valfull}). With both patches, every ML-DSA case behaves as the classical ones do, and a manifest timestamped by an ES256 authority validates in \NvalVerifyLo{} to \NvalVerifyHi{}\,ms on the M1 Pro. The fully post-quantum credential, an ML-DSA-65 signer timestamped by an ML-DSA-65 authority under an ML-DSA-65 anchor, validates as \emph{Trusted} with a trusted timestamp, in \NvalPqTsaMs\,ms, and it still does after the signer's certificate has expired, which is the job a timestamp exists to do. Through the signing patch alone the same timestamp is reported untrusted, because its authority's path cannot be checked.
The library rejects a modified image, a flipped signature bit, a mismatched algorithm and an expired leaf whose timestamp is missing, covers other data or comes from an untrusted authority, so a rejected timestamp never lends an expired signer historical validity; it stops short of \emph{Trusted} for a chain to an unknown anchor, and refuses to sign with a leaf that fails its own certificate profile. It
accepts the leaf alone when the leaf's issuer is the anchor, the
configuration Section~\ref{sec:carrying} recommends.

Two results depart from C2PA~2.5, for every algorithm. A leaf whose only
extended key usage is \texttt{id-kp-timeStamping} validates as a claim signer,
even under a policy that lists only \texttt{c2pa-kp-claimSigning}: the
library's key-usage check accepts e-mail protection, time-stamping and OCSP
signing unconditionally, which lets a time-stamping authority's certificate
sign manifests. In the other direction, a time-stamping authority whose
certificate lacks the time-stamping usage makes the whole manifest
\emph{Invalid}, where the specification says to report the timestamp as
untrusted and ignore it. Neither is specific to post-quantum algorithms: both appear unchanged in the same commit built without any patch, in the latest tagged release, c2pa-v0.91.9, and in the command-line tool c2patool~0.28.4. We reported the second publicly~\cite{c2pars3005}. The first lets a time-stamping key sign manifests that validate as \emph{Trusted}, so the released data withhold its signed file and certificate until it is fixed; minimal vectors for the second and raw status codes for both are released. A
manifest carrying a thumbprint, or no
credential, fails with a reader error rather than a validation status, and a
corrupt or truncated signature with an error about time-stamping; a
50-certificate chain is reported as untrusted, and a 1\,MB certificate validates in
\NvalBigCertMs\,ms using \NvalBigCertMB\,MB of memory.

Manifests the library writes with leaf and intermediate are larger than our
direct structures over the same image by \Nlibgap{} bytes at every ML-DSA
parameter set, and by \NlibgapClassLo{} to \NlibgapClassHi{} under the
classical algorithms (Table~\ref{tab:c2parust} in the supplement). Most of the difference is padding: c2pa-rs pads the signature, in its unprotected header, to the space the signer reserved before signing, and our harness reserves the signature and chain plus 2\,KiB, which leaves \NlibpadLo{} to \NlibpadHi{} bytes of padding;
its claim, assertions and framing take \Nliboutside{} bytes against our
\Nouroutside, and its test certificates, which lack the production extensions,
are slightly smaller than ours. The difference does not grow with the
signature, so every conclusion about how cost scales with the algorithm holds
for the library's output.

%% ---------------------------------------------------------------------------
\section{The cost of a post-quantum content credential}
\label{sec:cost}

\begin{table}[t]
\caption{Manifest cost by signature algorithm with the signer's certificate and its intermediate in \texttt{x5chain}, the root omitted as C2PA~2.5 says it should be. Sizes are exact; the percentage is the median over the Kodak set with its range. \emph{vs ES256} and \emph{vs PS256} compare with the two algorithms deployed signers use (Section~\ref{sec:prodchains}). Signing time is the signature over the COSE \texttt{Sig\_structure} on Apple M1 Pro, through liboqs for ML-DSA and OpenSSL for the classical rows. Times are medians over 10 independent runs, each the median of its own repetitions; Table~\ref{tab:ci} in the supplement gives their intervals across runs.}
\label{tab:pq}
\small
\setlength{\tabcolsep}{4pt}
\begin{tabular}{lrrrrrrrr}
\toprule
 & COSE & Sig. & Chain & Manifest & \% of image & vs & vs & Sign\\
Algorithm & id & (B) & (B) & (B) & median (range) & ES256 & PS256 & (ms)\\
\midrule
Ed25519 & -8 & 64 & 1,387 & 2,544 & 2.37 (1.35--3.29) & 1.0$\times$ & 0.7$\times$ & 0.030\\
ES256 & -7 & 64 & 1,501 & 2,658 & 2.48 (1.41--3.44) & 1.0$\times$ & 0.7$\times$ & 0.022\\
PS256 & -37 & 256 & 2,291 & 3,642 & 3.39 (1.94--4.71) & 1.4$\times$ & 1.0$\times$ & 0.535\\
ML-DSA-44 & -48 & 2,420 & 8,723 & 12,238 & 11.40 (6.51--15.83) & 4.6$\times$ & 3.4$\times$ & 0.074\\
ML-DSA-65 & -49 & 3,309 & 11,781 & 16,185 & 15.08 (8.61--20.93) & 6.1$\times$ & 4.4$\times$ & 0.116\\
ML-DSA-87 & -50 & 4,627 & 15,697 & 21,419 & 19.95 (11.39--27.70) & 8.1$\times$ & 5.9$\times$ & 0.165\\
\bottomrule
\end{tabular}
\end{table}


Table~\ref{tab:pq} is the central result, and it has to be read against what
signers use today, which is ES256 and, about half as often, RSA-2048 PS256
(Section~\ref{sec:prodchains}). Against ES256 an ML-DSA manifest is
\NvsESDSAfortyfour{}, \NvsESDSAsixtyfive{} and \NvsESDSAeightyseven{} times
larger; against PS256, \NvsPSDSAfortyfour{}, \NvsPSDSAsixtyfive{} and
\NvsPSDSAeightyseven{} times. Either way the manifest grows far more slowly than
the signature, which grows from 64 bytes under ES256 and 256 under RSA-2048 to
3,309 under ML-DSA-65.

The chain column says why. Leaf and intermediate together are
\NchainBytesDSAsixtyfive{} bytes under ML-DSA-65, \NchainShareDSAsixtyfive\% of the signed COSE structure and \NchainShareManDSAsixtyfive\% of the whole manifest once JUMBF framing and the claim are added, against a 3,309-byte signature; under ES256 the
same two certificates are \NchainShareES\% of a much smaller structure. Each
certificate carries a public key and its issuer's signature, so every
certificate in the chain repeats the growth that the manifest signature
shows once. Post-quantum cryptography did not create this imbalance; it scaled
an existing one until it became the cost.

Signing time is not part of the argument. Through liboqs, ML-DSA-44 signs in
\NsignMsDSAfortyfour\,ms, against \NsignMsEd\,ms for Ed25519 and
\NsignMsPS\,ms for PS256 through OpenSSL on the same machine.

\subsection{The denominator problem}
\label{sec:denominator}

A percentage of a Kodak image invites a division the reader will do anyway.
Kodak images are 768$\times$512 and a median of about 108\,kB; a phone
photograph is closer to 4\,MB, where the same ML-DSA-65 manifest that is
\NpctimgDSAsixtyfive\% of a Kodak image is under half a per cent. Once
thumbnails are disabled, manifest size does not depend on the pixels: it is
fixed by the algorithm, by how much of the chain travels, and by the
assertions the generator writes, which for a capture are a handful of hashes
and actions. Ingredients, edit histories and large assertions add bytes of
their own that no choice of algorithm changes. The percentage is therefore an
exact curve, and Figure~\ref{fig:sizeladder} plots it.

\begin{figure}[t]
\centering
\includegraphics[width=0.86\linewidth]{figures/sizeladder}
\Description{Log-log plot of manifest size as a percentage of file size
against file size from 10 kB to 20 MB. Falling straight lines for ML-DSA-65
with timestamp and OCSP, ML-DSA-65 timestamped, ML-DSA-65 with leaf and
intermediate, ML-DSA-65 leaf only, ML-DSA-65 thumbprint, and PS256 timestamped.
A shaded band marks the Kodak images; vertical guides mark a thumbnail, a
social re-encode, a 1080p JPEG and a phone original.}
\caption{Manifest cost against file size. Sizes are measured
(Tables~\ref{tab:timestamp} and~\ref{tab:carriage}) and do not depend on
image content once thumbnails are disabled, so each line is exact. The shaded
band is the Kodak set; the vertical guides are reference sizes, not
measurements. Both axes are logarithmic, so the vertical distance between two
lines is the ratio of their manifest sizes at every file size.}
\label{fig:sizeladder}
\end{figure}

A reader with phone originals in mind will find the migration affordable, and
for that case they are right. A platform storing derivatives, a camera signing
burst frames or a pipeline signing video operates at the left of the figure,
where the manifest is a double-digit percentage of the payload. The two
readings disagree about file size, not about cryptography.

\subsection{The manifest a deployment ships}
\label{sec:timestamp}

\begin{table}[t]
\caption{The manifest a deployment ships: signature, then RFC~3161 timestamp, then stapled revocation evidence, as most production signatures carry a timestamp and a quarter also carry revocation evidence (Section~\ref{sec:prodchains}). Signer and time-stamping authority each carry their certificate and intermediate, roots omitted; the \emph{+ OCSP} column adds a response from a delegated responder, with the responder's certificate inside it, the form the reference implementation honours. Every object is real and signed under the named algorithm. Figure~\ref{fig:compose} breaks the bytes down by role.}
\label{tab:timestamp}
\small
\setlength{\tabcolsep}{4pt}
\begin{tabular}{lrrrrrr}
\toprule
 & Signature & + timestamp & + OCSP & \multicolumn{2}{c}{Timestamped vs} & \% of\\
\cmidrule(lr){5-6}
Algorithm & (B) & (B) & (B) & ES256 & PS256 & image\\
\midrule
Ed25519 & 2,544 & 4,519 & 5,549 & 1.0$\times$ & 0.7$\times$ & 4.2\\
ES256 & 2,657 & 4,722 & 5,819 & 1.0$\times$ & 0.7$\times$ & 4.4\\
PS256 & 3,642 & 6,687 & 8,371 & 1.4$\times$ & 1.0$\times$ & 6.2\\
ML-DSA-44 & 12,238 & 23,915 & 30,977 & 5.1$\times$ & 3.6$\times$ & 22.2\\
ML-DSA-65 & 16,185 & 31,809 & 41,289 & 6.7$\times$ & 4.8$\times$ & 29.5\\
ML-DSA-87 & 21,419 & 42,277 & 55,033 & 9.0$\times$ & 6.3$\times$ & 39.3\\
\bottomrule
\end{tabular}
\end{table}


Deployed manifests carry more than one signature. \NcorpProdTs{} of the
\NcorpProdSigs{} production signatures we read carry an RFC~3161
timestamp~\cite{rfc3161}, which is a second signature with certificates of
its own, and the specification asks for one whenever revocation status is
attached~\cite{c2pacc24}. Revocation evidence is a third signature, a stapled
OCSP response~\cite{rfc6960}, and \NcorpProdOcsp{} of the production
signatures carry it; validators need not check it, the gap that Golaszewski et
al.\ identify~\cite{golaszewski2026}, which makes the last column of
Table~\ref{tab:timestamp} the configuration that closing the gap would make
universal. RFC~6960 lets the issuing CA sign the response itself, and then
it carries no certificate, since the verifier already has the CA's from
\texttt{x5chain}. C2PA~2.5 requires the certificate that signs OCSP
responses to hold the OCSP-signing usage, and the reference implementation
skips a CA-signed response (Table~\ref{tab:validation}), so the column uses a
delegated responder, which carries its own certificate:
\NocspDelDSAsixtyfive{} bytes under ML-DSA-65 against
\NocspCADSAsixtyfive{} for the CA-signed form.

The timestamp multiplies the certificate term rather than adding to it. Under
ES256 it takes the manifest from \NtsSigES{} to \NtsTstES{} bytes; under
ML-DSA-65, from \NtsSigDSAsixtyfive{} to \NtsTstDSAsixtyfive{}, a factor of
\NtsGrowthDSAsixtyfive{}, because the time-stamping authority's certificates
cost what the signer's do. Figure~\ref{fig:compose} breaks a timestamped
manifest down by role: under every post-quantum parameter set, certificates
are the largest part and the two manifest signatures a minority.

\begin{figure}[t]
\centering
\includegraphics[width=\linewidth]{figures/compose}
\Description{Horizontal stacked bars, one per algorithm from Ed25519 to
ML-DSA-87, showing the bytes of a timestamped manifest split into manifest
signature, signer leaf certificate, signer intermediate certificate, TSA
signature, TSA certificates, and structure. The bars grow from \NtsTstkBEd{} kB for
Ed25519 to \NtsTstkBDSAeightyseven{} kB for ML-DSA-87, and certificates are the largest
segments in the post-quantum bars.}
\caption{What a timestamped manifest is made of, by role: the \emph{+ timestamp} column of Table~\ref{tab:timestamp}, roots omitted, without the OCSP response. Structure covers COSE and JUMBF
wrapping, the claim, and the CMS and TSTInfo encoding of the timestamp.}
\label{fig:compose}
\end{figure}

\subsection{What signing actually costs}
\label{sec:hash}

A content credential is not only a signature. The claim contains a hard
binding, a hash over the content, computed before anything is signed. On a
4\,MB phone photograph SHA-256 takes \NhashPhone\,ms on the M1 Pro, so the hash
is \NhashSharePhone\% of hashing and ML-DSA-65 signing together; on a Kodak
image it is \NhashShareKodak\% (Table~\ref{tab:hash} in the supplement). The
platform has hardware SHA-256 and no hardware support for lattice arithmetic,
the configuration least favourable to this conclusion. On images large enough
to be worth signing, the post-quantum signature costs less than the hash the
format already requires; the signer's cost of migrating is bytes, and time
only at small sizes and high rates, which Section~\ref{sec:edge} measures.

\subsection{Beyond ML-DSA}
\label{sec:algsweep}

\begin{table}[t]
\caption{Signature schemes at hierarchy depth three, leaf and intermediate carried. \emph{Chain} is the share of the signed structure taken by certificates; \emph{Embed.} and \emph{Ref.} are the manifest with the chain and with a thumbprint, as a percentage of the median image. Every certificate is real; Falcon carries a placeholder object identifier and its padded signature length. A dash means no COSE identifier is assigned. Signing times through liboqs on Apple M1 Pro, over the claim alone (Section~\ref{sec:timing}). Times are medians over 10 independent runs; their intervals across runs are in the released data.}
\label{tab:algsweep}
\small
\setlength{\tabcolsep}{4pt}
\begin{tabular}{llrrrrrr}
\toprule
 & COSE & Sig. & Chain & Chain & Embed. & Ref. & Sign\\
Algorithm & id & (B) & (B) & (\%) & (\% img) & (\% img) & (ms)\\
\midrule
Ed25519 & -8 & 64 & 1,387 & 94.2 & 2.4 & 1.1 & --\\
ES256 & -7 & 64 & 1,499 & 94.6 & 2.5 & 1.1 & --\\
PS256 & -37 & 256 & 2,291 & 89.1 & 3.4 & 1.3 & --\\
\midrule
ML-DSA-44 & -48 & 2,420 & 8,723 & 78.1 & 11.4 & 3.3 & 0.061\\
ML-DSA-65 & -49 & 3,309 & 11,781 & 78.0 & 15.0 & 4.1 & 0.096\\
ML-DSA-87 & -50 & 4,627 & 15,697 & 77.2 & 19.9 & 5.3 & 0.138\\
Falcon-512 & -- & 666 & 4,372 & 86.4 & 5.7 & 1.7 & 0.155\\
Falcon-1024 & -- & 1,280 & 7,398 & 85.0 & 9.1 & 2.2 & 0.304\\
SLH-DSA-SHA2-128f & -- & 17,088 & 35,558 & 67.5 & 49.9 & 16.9 & 21\\
SLH-DSA-SHA2-128s & -- & 7,856 & 17,094 & 68.5 & 24.2 & 8.3 & 429\\
SLH-DSA-SHA2-192s & -- & 16,224 & 33,862 & 67.6 & 47.5 & 16.1 & 757\\
SLH-DSA-SHA2-256s & -- & 29,792 & 61,030 & 67.2 & 85.4 & 28.7 & 672\\
\bottomrule
\end{tabular}
\end{table}


Table~\ref{tab:algsweep} holds the chain share between \NsweepShareLo\% and
\NsweepShareHi\% for signatures whose sizes span a factor of
\NsweepSigFactor. SLH-DSA is unsuitable for per-image provenance at the
burst and frame rates this paper considers: SHA2-256s produces a manifest of
\NslhBigPct\% of the median image, and the compact variants sign in hundreds
of milliseconds; a workflow that signs a few documents an hour could still
use it. Falcon-512 is close to classical cost, a manifest of \NfalconEmbPct\%
of the median image with the chain embedded, but FN-DSA is unpublished, its
arithmetic is being changed from floating to fixed point, and the timings here
are for the submitted algorithm~\cite{nist_fips206_status,nist_fips206_plan}.

Two controls a reader will ask for are in the supplement. Composite
signatures, built from their two components with real certificates
(Table~\ref{tab:composite}), cost the sum of their parts: an
ML-DSA-65 and ES256 composite is \NcompositeCose{} bytes of COSE against
\NcoseDSAsixtyfive{} for ML-DSA-65 alone, so hedging against a lattice
break costs a few per cent of the manifest. General-purpose compression
(Table~\ref{tab:compress}) removes \NcompressDSAsixtyfive\% of an ML-DSA-65
signature structure and \NcompressEd\% of an Ed25519 one: what compresses is
X.509 names and structure, and post-quantum keys and signatures do not.

\subsection{Chain depth}
\label{sec:chaindepth}

Each further level of the hierarchy adds one certificate to \texttt{x5chain}
and costs the same again (Table~\ref{tab:chaindepth} in the supplement).
Under ML-DSA-65 the difference between a two-level and a four-level hierarchy
is \NdepthSpread{} percentage points of a Kodak image, the order of the
difference between ML-DSA-44 and ML-DSA-87. Leaving a root out of
\texttt{x5chain} that the verifier already holds changes no trust decision;
shortening the hierarchy itself is a change of PKI design, and its security
consequences belong to the operator rather than to this measurement.

\subsection{The chains deployments actually ship}
\label{sec:prodchains}

The configuration costed so far is ours, and a reader is entitled to ask what
signers actually ship. We read the signatures out of \NcorpProdAssets{} files
signed by production software, \NcorpProdGens{} signing applications from
camera and phone makers, AI image services, editing tools and news
publishers, collected from Wikimedia Commons, the IPTC's media-provenance samples, public datasets and vendor posts, and out of the C2PA public test files. This is a convenience sample of what is publicly posted, not a census of signers. Every signature in a file is counted, including those of ingredients, because the file carries them. A chain is identified by its issuer, the leaf's issuer name and every certificate above the leaf, so a signer that issues a leaf per image counts once. Table~\ref{tab:corpus} gives the result by the application that wrote the file; a chain that reaches files of several groups, such as a Pixel camera chain in an image later edited with Microsoft's tools, appears in each row, and the \emph{Distinct} row counts every chain once; \NcorpMultiGroup{} of the \NcorpProdChains{} chains appear in more than one group. The supplement lists every chain and where the files came from (Section~\ref{sec:supp-corpus}); the released inventory gives every file's source URL, licence and SHA-256.

\begin{table*}[t]
\caption{Certificate chains in 317 production-signed files from 28 signing applications, and in the C2PA public test files, read against the C2PA trust list of 29~October 2026, a convenience sample of public files (Section~\ref{sec:supp-corpus} of the supplement lists the chains and sources, and the released inventory every file with its URL, licence and SHA-256). A chain is identified by its issuer: the leaf's issuer name and every certificate after the leaf, so leaves issued per image under one CA count once. Every signature in a file counts, the active manifest's and its ingredients', under the group of the application that wrote the file; a chain carried by files of several groups (Google's camera chains in files that Microsoft's tools wrote, for example) is counted in each, so the rows do not add up, and the \emph{Distinct} row counts each chain once. \emph{Sent}: certificates in \texttt{x5chain}. \emph{Root}: chains that carry a self-signed root. \emph{Listed}: chains whose path ends at a trust-list anchor. \emph{Issuer}: chains whose leaf was issued directly by an anchor, the condition for conformant leaf-only carriage. \emph{2.x}: files with a version-2 claim. \emph{33p}: signatures carrying \texttt{x5chain} under label~33 in the protected bucket. \emph{TST}: signatures with an RFC~3161 token, and its median size. \emph{OCSP}: signatures carrying revocation evidence.}
\label{tab:corpus}
\footnotesize
\setlength{\tabcolsep}{3.1pt}
\begin{tabular}{lrrrrlrrrrrrrr}
\toprule
 & & & \multicolumn{2}{c}{Chains} & & & & & & & \multicolumn{2}{c}{TST} & \\
\cmidrule(lr){4-5}\cmidrule(lr){12-13}
Signer & Files & Sigs & ES256 & PS256 & Sent & Root & Listed & Issuer & 2.x & 33p & n & B & OCSP\\
\midrule
Google & 105 & 141 & 5 & 0 & 2 & 0 & 5 & 2 & 105 & 141 & 141 & 2,015 & 43\\
OpenAI & 53 & 83 & 3 & 1 & 2, 3 & 1 & 3 & 0 & 51 & 83 & 31 & 5,259 & 23\\
Microsoft & 66 & 94 & 6 & 1 & 1, 2, 3 & 1 & 4 & 2 & 6 & 27 & 85 & 5,931 & 20\\
Adobe & 26 & 27 & 0 & 2 & 2 & 0 & 0 & 0 & 0 & 27 & 27 & 3,664 & 26\\
News publishers & 24 & 35 & 4 & 6 & 1, 2, 3 & 2 & 0 & 0 & 3 & 32 & 25 & 5,625 & 2\\
Other & 43 & 45 & 4 & 1 & 1, 2 & 0 & 1 & 0 & 3 & 45 & 43 & 5,828 & 0\\
\midrule
\emph{Distinct} & 317 & 425 & 15 & 8 & 1, 2, 3 & 4 & 9 & 2 & 168 & 355 & 352 & 3,663 & 114\\
\midrule
C2PA test files & 25 & 41 & 1 & 3 & 2, 3 & 1 & 0 & 0 & 0 & 0 & 39 & 5,940 & 0\\
\bottomrule
\end{tabular}
\end{table*}


\emph{The test files are not production output.} The C2PA public test files
hold \NcorpTestChains{}
distinct chains, almost all Adobe test certificates, every file has a
version-1 claim, and none reaches the current trust list. They show what a
conformance test looks like, not what signers emit.

\emph{ECDSA leads, and RSA is far from gone.} Of the
\NcorpProdChains{} distinct production chains, \NcorpProdEs{} sign with ES256
and \NcorpProdPs{} with PS256, and \NcorpProdSigEs{} of the
\NcorpProdSigs{} signatures are ES256. Google's phones and media services use
ES256 and OpenAI's image service uses both; RSA, mostly 2048-bit, persists in
Adobe's tools, Microsoft's image service and most newsrooms. Both are
therefore the baseline, and Table~\ref{tab:pq} reports growth against each.

\emph{Depth is three, and the root mostly stays out.} Signers carry the leaf and one intermediate in \NcorpProdTwoCert{} of the \NcorpProdChains{} chains; \NcorpProdRoot{} of the
\NcorpProdChains{} still carry the root, which the specification says to
leave out, and \NcorpProdOneCert{} carry the leaf alone, none of them with an
issuer on the trust list, so no validator using that list can take them to
\emph{Trusted}. The median chain carried is \NcorpProdChainMed{}
bytes.

\emph{The leaf-only condition already holds for some signers.}
Of the \NcorpProdChains{} chains, \NcorpProdListed{} end at an anchor on the
C2PA trust list, which was established in 2025 and does not yet cover many
earlier signers. Two of them, both Google's, are issued directly by a listed
anchor, so the path a validator checks is already one signature; they also
appear in the Microsoft row because Pixel images were edited there. Both still
carry that anchor beside the leaf, which the specification says not to
include, and dropping it would make them leaf-only and save about 730 bytes
each.

\emph{Timestamps and revocation are the ordinary case.} Of the
\NcorpProdSigs{} production signatures, \NcorpProdTs{} carry an RFC~3161 token, of
\NcorpProdTstLo{} to \NcorpProdTstHi{} bytes with a median of
\NcorpProdTstMed{}: about 2\,kB for an ECDSA time-stamping authority with two
certificates, as on Google's devices, and about 6\,kB for an RSA-4096 one with
three. Our ES256 and PS256 tokens, \NtokenBytesES{} and \NtokenBytesPS{}
bytes, sit at the small end of that range. \NcorpProdOcsp{} production
signatures each staple at least one OCSP response, so the OCSP column of
Table~\ref{tab:timestamp} describes what a quarter of today's signatures
already carry, not a hardened variant.

\emph{The header has moved to the RFC~9360 form.} Of the
\NcorpProdNxfive{} production signatures, \NcorpProdPthree{} carry \texttt{x5chain} under the
integer label~33 in the protected bucket, as C2PA~2.5 says generators shall;
the rest use the deprecated string label in the unprotected bucket, all but
\NcorpProdVtwoString{} of them in version-1 claims, which validators must
still accept.

\emph{What ML-DSA-65 would add to these files.} Because size is additive in
the components (Equation~\ref{eq:size}), the migration can be priced on the
files themselves. For each of the \NcpqFiles{} production files we kept
everything it carries, claims, assertions, ingredients, thumbnails and the
number of certificates each signer and each time-stamping authority sends, and
replaced every signature, carried certificate and timestamp token with our
measured ML-DSA-65 counterpart of the same role, so production names and
extensions give way to our profile. Three things make this an estimate rather
than a measurement of re-signed files, all of which widen it: the files are not
actually re-signed; a production certificate may differ from our profile by a
few hundred bytes, which enters the per-certificate total directly
(Section~\ref{sec:limitations}); and stapled OCSP responses, whose size the
files do not expose, are left classical, with a post-quantum upper bound given
separately (Table~\ref{tab:corpusstrata}). Percentages here are of the original
file. The median file is \NcpqFileMedMB\,MB; the median growth is
\NcpqAddMedkB\,kB, or \NcpqPctMed\% of a file, and the median
manifest store grows by a factor of \NcpqStoreGrowth. The median hides the
tail: files with several signatures, such as edited images that keep their
ingredients' credentials, grow by up to \NcpqAddMaxkB\,kB, and
\NcpqOverTen{} of the \NcpqFiles{} files would grow by more than a tenth
(90th percentile \NcpqPctPNinety\%). Of those, \NcpqOverTenTop{} are generated
images of about 200\,kB from one AI image service; among the other
\NcpqFilesOther{} files, \NcpqOverTenOther{} grow that much. Size decides the share: files under
500\,kB grow by a median \NcpqSmallMedPct\%, files over 5\,MB by
\NcpqBigMedPct\%, and the \NcpqThreeSigFiles{} files that keep three or more
signatures gain a median \NcpqThreeSigMedkB\,kB (Table~\ref{tab:corpusstrata}
and Figure~\ref{fig:corpusgrowth} in the supplement). Giving the
\NcpqOcspSigs{} stapled responses a full post-quantum delegated response
each, with nothing subtracted for the classical one, raises the median growth
of the \NcpqOcspFiles{} files that carry one from \NcpqOcspFilesMedkB{} to
\NcpqOcspFilesBoundMedkB\,kB, an upper bound. For the services that sign small
images at volume, the post-quantum credential is a double-digit share of every
file they serve. The released data list every file.

%% ---------------------------------------------------------------------------
\section{Carrying the signer}
\label{sec:carrying}

If certificates carry the cost, the question is how many of them have to
travel. Table~\ref{tab:carriage} measures four answers under two placements
of the trust anchor.

\begin{table*}[t]
\caption{Carrying the signer. Manifest bytes by carriage form and algorithm, three-certificate hierarchy. \emph{Path} is the number of signatures path validation checks, set by where the trust anchor sits; \emph{Holds} is what the verifier must already have; \emph{int.} is the intermediate. $^a$C2PA~2.5 says the root should not be included. $^b$A reading the specification does not exclude, accepted by the reference implementation when the leaf's issuer is a configured anchor; clarification requested~\cite{c2paspec142}. Otherwise the leaf-only form does not conform.}
\label{tab:carriage}
\footnotesize
\setlength{\tabcolsep}{3.5pt}
\begin{tabular}{llclrrrrrr}
\toprule
 & & & &  &  &  & \multicolumn{3}{c}{ML-DSA}\\\cmidrule(l){8-10}
\texttt{x5chain} & Holds & Path & Conf. & Ed25519 & ES256 & PS256 & 44 & 65 & 87\\
\midrule
Leaf, int., root & root & 2 & yes$^a$ & 2,981 & 3,155 & 4,535 & 16,339 & 21,815 & 29,007\\
Leaf, int. & root & 2 & yes & 2,544 & 2,658 & 3,642 & 12,238 & 16,185 & 21,419\\
Leaf (issuer is anchor) & issuing CA & 1 & yes$^b$ & 1,849 & 1,906 & 2,495 & 7,875 & 10,293 & 13,569\\
Leaf (int.\ provisioned) & root, int. & 2 & no & 1,849 & 1,906 & 2,495 & 7,875 & 10,293 & 13,569\\
Thumbprint & root, int., leaf & 2 & no & 1,185 & 1,185 & 1,379 & 3,543 & 4,432 & 5,750\\
\midrule
\multicolumn{4}{l}{Saving, leaf only (\%)} & 27 & 28 & 31 & 36 & 36 & 37\\
\multicolumn{4}{l}{Saving, thumbprint (\%)} & 53 & 55 & 62 & 71 & 73 & 73\\
\multicolumn{4}{l}{Cost of the root (\%)} & 17 & 19 & 25 & 34 & 35 & 35\\
\bottomrule
\end{tabular}
\end{table*}


\paragraph{Leaf, intermediate and root} What the Adobe chains in the public
test corpus carry. The specification says the anchor's certificate should not
be included, and it costs \NcostRootDSAsixtyfive\% more than leaving it out
under ML-DSA-65 while changing no trust decision, since the verifier must
hold the root anyway. Dropping it is free.

\paragraph{Leaf and intermediate} The form C2PA~2.5 requires and the baseline
of every comparison in this paper.

\paragraph{Leaf only} The specification requires every intermediate CA certificate and says the trust anchor's should be left out. When the CA that issued the
leaf is itself on the trust list, it is the anchor, there is no intermediate
to carry, and the leaf alone is the conformant form. The C2PA trust list is
built to allow this: \NtlIssuing{} of its \NtlAnchors{} anchors are issuing
CAs rather than roots, among them Google's capture and media CAs, DigiCert's
claim-signing CAs and Adobe's product CA~\cite{c2patrustlist}. In that case
the manifest is \NsaveLeafDSAsixtyfive\% smaller under ML-DSA-65, and path
validation checks one signature instead of two, because an anchor's own
signature is never checked. When the leaf's issuer is not an anchor, carrying
the leaf alone works only for a verifier that already holds the intermediate,
which no specification provides for; Table~\ref{tab:validation} shows the
reference implementation returning \emph{Valid} rather than \emph{Trusted} in
exactly that case.

\paragraph{Thumbprint} The \texttt{x5t} header and no certificate, a further
\NsaveThumbDSAsixtyfive\% saving against leaf and intermediate. It is not
conformant, and the reference implementation reads it as a missing
credential.

Section~\ref{sec:prodchains} measures how often the conformant leaf-only
condition holds for production signers today: \NcorpProdListed{} of the
\NcorpProdChains{} distinct production chains we found end at a trust-list
anchor, and for \NcorpProdIssuer{} of them the leaf's issuer is the anchor. Our recommendation follows from both
measurements. Generators should never carry the root. A CA that issues
post-quantum signing certificates should have its issuing CA listed as a
trust anchor, as several already do for classical ones, and generators under
such a CA should carry the leaf alone; the saving is a third of the manifest
and one of two path signatures, and it needs no change to any specification.

\subsection{What each option gives up}
\label{sec:tradeoffs}

\emph{The trust list does not resolve a thumbprint.} It holds anchors, not
signers' leaf certificates, so a verifier holding only an \texttt{x5t} has
nowhere specification-defined to fetch the leaf from, and some signers issue
a leaf per image~\cite{google2025pixel}, for which there could be no such
directory.

\emph{Offline verification of an unknown signer.} A verifier with no network
that has never seen the signer can verify leaf and intermediate, can verify
the leaf alone when the leaf's issuer is an anchor it holds, and cannot verify
a thumbprint. Archival readers opening a file decades later are the case where
provenance matters most and fetching is least possible.

\emph{Privacy.} Remote thumbprint resolution reports every verification to
whoever hosts the resolver: which verifier looked up which signer, and when.
For a journalist checking a document, or a reader of politically sensitive
media, that is a disclosure the embedded forms do not make.

\subsection{Where verification time goes}
\label{sec:verifytime}

\begin{table*}[t]
\caption{Where the cryptographic work of validating one timestamped image goes, in milliseconds on Apple M1 Pro. \emph{Path} checks leaf and intermediate up to the root anchor; \emph{TST} and \emph{TSA path} are the timestamp's signature and certificate path; \emph{Hash} is SHA-256 over the content, for the 13th of the 24 Kodak images by size (113,540\,B) and for a 4\,MiB phone JPEG. Assertion hashing, ingredients, revocation and policy checks are outside the boundary. \emph{Total} is the sum of the row for each image. Times are medians over 10 independent runs, each the median of its own repetitions; Table~\ref{tab:ci} in the supplement gives their intervals across runs.}
\label{tab:verify}
\small
\setlength{\tabcolsep}{4pt}
\begin{tabular}{lrrrrrrrrr}
\toprule
 & Decode & Path & Manifest & TST & TSA & \multicolumn{2}{c}{Hash} & \multicolumn{2}{c}{Total}\\
\cmidrule(lr){7-8}\cmidrule(lr){9-10}
Algorithm & & & signature & & path & median & 4\,MiB & median & 4\,MiB\\
\midrule
Ed25519 & 0.001 & 0.181 & 0.081 & 0.080 & 0.180 & 0.046 & 1.70 & 0.571 & 2.23\\
ES256 & 0.001 & 0.148 & 0.056 & 0.054 & 0.148 & 0.046 & 1.71 & 0.456 & 2.12\\
PS256 & 0.001 & 0.069 & 0.021 & 0.016 & 0.069 & 0.046 & 1.71 & 0.223 & 1.89\\
ML-DSA-44 & 0.002 & 0.170 & 0.043 & 0.070 & 0.171 & 0.046 & 1.71 & 0.501 & 2.16\\
ML-DSA-65 & 0.002 & 0.251 & 0.068 & 0.108 & 0.251 & 0.046 & 1.70 & 0.729 & 2.39\\
ML-DSA-87 & 0.003 & 0.380 & 0.110 & 0.171 & 0.380 & 0.046 & 1.71 & 1.089 & 2.76\\
\bottomrule
\end{tabular}
\end{table*}


Certificates dominate verification for the same reason they dominate bytes,
but how much depends on what is counted. For the manifest signature and the
signer's path alone, the measurement most comparable to a TLS handshake, the
path is \NvPathShareSigDSAsixtyfive\% of the ML-DSA-65 work in
Table~\ref{tab:verify}, against the two thirds that two path signatures and
one manifest signature would suggest; the excess is that the path goes
through OpenSSL's X.509 stack, whose ML-DSA verification is slower than
liboqs's, and a verifier cannot choose a faster one for certificates because
liboqs has no ASN.1 layer. That fraction is conditional on the hierarchy: with
the issuing CA as anchor the path is one signature, and with a cached signer
(Section~\ref{sec:cache}) it is none.

Complete validation of a timestamped image adds the timestamp's signature and
path and the content hash, and Figure~\ref{fig:latency} shows the result. For a mid-sized Kodak image, the 13th of 24 by size, the signer's path falls to \NvPathShareDSAsixtyfive\% of
the cryptographic work under ML-DSA-65; for a 4\,MiB phone JPEG, hashing the
content dominates and the algorithm barely shows, \NvCompletePhonePS\,ms
under PS256 against \NvCompletePhoneDSAsixtyfive\,ms under ML-DSA-65. The
reference implementation, which also parses JUMBF and checks every assertion,
takes \NvalVerifyLo{} to \NvalVerifyHi\,ms per timestamped image
(Section~\ref{sec:patch}).

\begin{figure}[t]
\centering
\includegraphics[width=\linewidth]{figures/latency}
\Description{Horizontal stacked bars, one per algorithm, showing the time to
validate one timestamped image split into decode, signer's path, manifest
signature, timestamp signature, TSA path and content hash for a mid-sized Kodak
image. Totals run from about 0.2 to 1.1 milliseconds, and the signer's path and the TSA path are the largest segments.}
\caption{The cryptographic work of validating one timestamped mid-sized Kodak
image (the 13th of 24 by size) on the M1 Pro (Table~\ref{tab:verify}). Assertion hashing, ingredients,
revocation and policy checks are outside the boundary.}
\label{fig:latency}
\end{figure}

\subsection{Caching a validated signer}
\label{sec:cache}

A verifier that sees many manifests can cache certificates it has already
validated and skip path validation for them. Whether that pays depends on how
signers use certificates, so Table~\ref{tab:cache} in the supplement runs a
verifier over a stream from \Ncachepublishers{} publishers, each with its own
issuing CA, under three leaf lifecycles. With one long-lived leaf per
publisher, a cache of 4,096 certificates makes ML-DSA-65 verification
\NcacheLongDSAsixtyfive{} times faster. With a new leaf every hundred manifests the gain is about the same, \NcacheRotDSAsixtyfive{} times. With a new leaf for
every image, as signers that issue a certificate per capture do, only the
issuing CA is ever found in the cache, and the gain falls to
\NcachePerDSAsixtyfive{} times.

A cache also has to forget. Half-way through each stream we removed one
publisher's issuing CA from trust. A cache whose entries carry the trust-list
snapshot they were validated under, as C2PA~2.5 now makes possible, rejected
every later manifest from that publisher. A cache without that tag went on
accepting \NnaiveLongDSAsixtyfive{} of the \NnaiveTotLongDSAsixtyfive{} that
followed,
because their certificates had been validated before the change. A snapshot tag covers only changes to the trust list. A cached validation also goes stale when a certificate on the path expires or is revoked while the list stays the same, so we ran the same stream over thirty simulated days with the most frequent publisher's signing certificate expiring, or revoked, half-way (Table~\ref{tab:cachestale} in the supplement). A cache keyed on the certificate alone accepted \NstaleExpNaive{} of the \NstaleExpTot{} manifests signed after expiry and \NstaleRevNaive{} of \NstaleRevTot{} after revocation; checking each entry against its path's notAfter stops the first and not the second. What works is to key an entry on the certificate, the trust-list snapshot and the validation policy, and to let it live only while every certificate on its path is valid and the revocation status it was validated with is current. With seven-day OCSP responses such a cache still accepted \NstaleRevBoundSeven{} manifests, the last \NstaleRevBoundSevenH{} hours after revocation, and with one-day responses \NstaleRevBoundOne{}, the last after \NstaleRevBoundOneH{} hours, which is the window any verifier that trusts a response until its nextUpdate has. A verifier that validates at a historical time, from a trusted timestamp, needs that time in the key as well.

\subsection{Verification across cores}
\label{sec:scale}

A platform checking credentials on upload sees manifests from many publishers
and asks how many cores that takes. We verify a stream from 500 publishers,
drawn from a Zipf distribution, so that a handful of busy publishers send a large share and most send a few each, on every core of the M1 Pro,
the Xeon VM and the Pi~5, with and without a cached signer
(Table~\ref{tab:scale} and Figure~\ref{fig:scale}).

\begin{table}[t]
\caption{Manifests verified per second, in thousands, on one core and on every core of each machine. A stream from 500 publishers, drawn Zipf(1), each with a three-certificate chain. \emph{Uncached} validates the embedded chain for every manifest; \emph{cached} validates each publisher's chain once per worker and then checks only the signature, through OpenSSL or, for ML-DSA, liboqs. x5chain carries leaf and intermediate; the root comes from the verifier's store. Medians of three independent runs of 20,000 verifications each.}
\label{tab:scale}
\begin{tabular}{lrrrrrr}
\toprule
 & \multicolumn{2}{c}{Uncached} & \multicolumn{2}{c}{Cached} & \multicolumn{2}{c}{Cached, liboqs}\\
\cmidrule(lr){2-3}\cmidrule(lr){4-5}\cmidrule(lr){6-7}
Algorithm & 1 core & all & 1 core & all & 1 core & all\\
\midrule
\multicolumn{7}{l}{\emph{Apple M1 Pro, all = 8 cores}}\\
Ed25519 & 3.6 & 21.1 & 11.2 & 54.6 & -- & --\\
ES256 & 4.4 & 25.7 & 15.4 & 73.1 & -- & --\\
PS256 & 9.5 & 55.6 & 40.4 & 171.3 & -- & --\\
ML-DSA-44 & 3.8 & 21.9 & 11.2 & 53.3 & 19.6 & 87.4\\
ML-DSA-65 & 2.5 & 14.7 & 7.4 & 36.1 & 12.6 & 56.9\\
ML-DSA-87 & 1.6 & 9.6 & 4.8 & 23.1 & 7.9 & 35.6\\
\midrule
\multicolumn{7}{l}{\emph{Intel Xeon VM, all = 2 cores}}\\
Ed25519 & 2.4 & 4.6 & 7.4 & 13.9 & -- & --\\
ES256 & 2.7 & 5.2 & 10.1 & 18.5 & -- & --\\
PS256 & 6.4 & 12.9 & 26.6 & 47.8 & -- & --\\
ML-DSA-44 & 2.1 & 4.2 & 6.5 & 12.4 & 20.0 & 32.6\\
ML-DSA-65 & 1.5 & 2.9 & 4.2 & 7.9 & 13.2 & 24.1\\
ML-DSA-87 & 1.0 & 1.9 & 2.8 & 5.3 & 9.8 & 16.6\\
\midrule
\multicolumn{7}{l}{\emph{Raspberry Pi 5, all = 4 cores}}\\
Ed25519 & 2.1 & 8.3 & 6.5 & 23.1 & -- & --\\
ES256 & 1.3 & 5.3 & 4.6 & 15.9 & -- & --\\
PS256 & 2.2 & 8.9 & 7.9 & 27.0 & -- & --\\
ML-DSA-44 & 1.3 & 5.1 & 3.9 & 13.7 & 8.5 & 25.8\\
ML-DSA-65 & 0.9 & 3.5 & 2.6 & 9.3 & 5.6 & 17.5\\
ML-DSA-87 & 0.6 & 2.3 & 1.7 & 6.1 & 3.6 & 11.2\\
\bottomrule
\end{tabular}
\end{table}


\begin{figure}[t]
\centering
\includegraphics[width=\linewidth]{figures/scale}
\Description{Two line charts of ML-DSA-65 verifications per second against
worker processes, uncached and cached, with one line each for the M1 Pro, the
Xeon VM and the Raspberry Pi 5. Throughput rises close to linearly on the
Pi's four cores and the M1 Pro's six performance cores.}
\caption{ML-DSA-65 verification throughput against worker processes on three
platforms, without and with a cached signer (Table~\ref{tab:scale}).}
\label{fig:scale}
\end{figure}

Without a cache, one core verifies \NscMUncDSAsixtyfive{} ML-DSA-65 manifests
a second on the M1 Pro, \NscXUncDSAsixtyfive{} on the Xeon VM and
\NscPUncDSAsixtyfive{} on the Pi~5. On the M1 Pro that is \NscMVsES\% of the
ES256 rate and \NscMVsPS\% of the PS256 rate; the Xeon (\NscXVsES\% and
\NscXVsPS\%) and the Pi (\NscPVsES\% and \NscPVsPS\%) give the same order.
The work parallelises:
\NscPMax{} workers on the Pi give \NscPSpeedUncDSAsixtyfive{} times one
worker's rate, and \NscMMax{} on the M1 Pro give \NscMSpeedUncDSAsixtyfive{}
times, its two efficiency cores adding less than a performance core. At one
core's rate, ten thousand ML-DSA-65 manifests a second need
\NscMCoresUncDSAsixtyfive{} M1 Pro cores, \NscXCoresUncDSAsixtyfive{} Xeon
vCPUs or \NscPCoresUncDSAsixtyfive{} Pi cores, against
\NscMCoresUncES{}, \NscXCoresUncES{} and \NscPCoresUncES{} for ES256. A
cached signer leaves only the manifest signature, \NscMCacDSAsixtyfive{} a
second per M1 Pro core through OpenSSL and \NscMOqsDSAsixtyfive{} through
liboqs, and cuts those counts to \NscMCoresOqsDSAsixtyfive,
\NscXCoresOqsDSAsixtyfive{} and \NscPCoresOqsDSAsixtyfive. These are throughputs for the manifest signature and the signer's path, the boundary of Table~\ref{tab:scale}. A complete validation of a timestamped median image, with the timestamp's signature and path and the content hash, is \NvCompleteFactorDSAsixtyfive{} times that work under ML-DSA-65 (Table~\ref{tab:verify}), and parsing, assertion checks and revocation lookups come on top; the core counts scale accordingly and are a lower bound for a full ingest pipeline. Within that boundary, post-quantum verification at upload scale is a matter of a few more cores, not of new hardware.

%% ---------------------------------------------------------------------------
\section{Signing and verifying where resources are short}
\label{sec:edge}

A camera, a body-worn recorder or a field encoder signs on something closer to
a Raspberry Pi~5 than to a laptop, and a doorbell or a sensor on something
smaller still. This section measures both. The Pi is an application processor,
not a secure element, and Section~\ref{sec:constrained} says what that leaves
open.

\subsection{Signing an encoded stream}
\label{sec:edge-video}

Cameras and encoders bind the encoded bitstream, not raw pixels, and C2PA
binds video through hashes of ISO BMFF segments rather than per-frame
signatures. C2PA~2.5 gives two methods for a live stream~\cite{c2pa25}. In the
first, every segment carries a manifest of its own, with a BMFF hash of the
segment, a live-video assertion carrying a sequence number and stream
identifier, and a signature with the signer's chain. In the second, the
stream's first manifest carries the chain and a session key bound to the
signer, and each segment carries only a small Verifiable Segment Info map,
its sequence number, hash and manifest identifier, signed by the session key.
We signed real streams both ways on one core of the Pi~5: \NstreamSeconds{} seconds of 1080p30 H.264 at 12\,Mbit/s and of 4K30 HEVC at 40\,Mbit/s, in fragmented MP4 with one fragment per second, so that a segment of $k$ seconds is $k$ fragments and even six-second segments give \NstreamSixSegs{} per pass. As an upper bound we also signed every encoded frame with the
chain attached, which C2PA does not define.

\begin{table*}[t]
\caption{Signing an encoded stream on the Raspberry Pi 5, one core. \emph{Frame} signs every encoded access unit with the chain attached (an upper bound, not a C2PA method); \emph{manifest} is C2PA~2.5 Section~15.4, a manifest per segment; \emph{VSI} is Section~15.5, the chain once and a session-key signature per segment. Provenance rate is the bitrate the signatures add; \emph{latency} is from segment close to signature (last fragment's hash plus signing). A player can authenticate a segment only once it holds all of it, so verification lags by at least the segment length. The provenance rate depends on the size of each signature and on how often one is sent, not on the video, so it is the same for both streams; the latency is not. Streams of 288\,s; medians over every segment of three passes.}
\label{tab:stream}
\small
\setlength{\tabcolsep}{3.2pt}
\begin{tabular}{llrrrrrrrrrrr}
\toprule
 & & Seg. & \multicolumn{2}{c}{Ed25519} & \multicolumn{2}{c}{PS256} & \multicolumn{2}{c}{ML-DSA-44} & \multicolumn{2}{c}{ML-DSA-65} & \multicolumn{2}{c}{ML-DSA-87}\\
\cmidrule(lr){4-5}\cmidrule(lr){6-7}\cmidrule(lr){8-9}\cmidrule(lr){10-11}\cmidrule(lr){12-13}
Stream & Method & (s) & kbit/s & ms & kbit/s & ms & kbit/s & ms & kbit/s & ms & kbit/s & ms\\
\midrule
1080p H.264 & frame & -- & 413.3 & 0.10 & 676.8 & 3.59 & 2,739.8 & 0.26 & 3,687.1 & 0.37 & 4,943.3 & 0.48\\
 & manifest & 1 & 22.3 & 1.07 & 31.1 & 4.56 & 99.9 & 1.22 & 131.5 & 1.32 & 173.3 & 1.44\\
 & manifest & 2 & 11.2 & 1.08 & 15.6 & 4.58 & 50.0 & 1.24 & 65.7 & 1.35 & 86.7 & 1.44\\
 & manifest & 4 & 5.6 & 1.06 & 7.8 & 4.56 & 25.0 & 1.22 & 32.9 & 1.33 & 43.3 & 1.44\\
 & VSI & 1 & 2.3 & 1.05 & 3.8 & 4.56 & 21.1 & 1.18 & 28.3 & 1.28 & 38.8 & 1.37\\
 & VSI & 2 & 1.1 & 1.06 & 1.9 & 4.57 & 10.6 & 1.20 & 14.1 & 1.29 & 19.4 & 1.39\\
 & VSI & 4 & 0.6 & 1.04 & 1.0 & 4.55 & 5.3 & 1.18 & 7.1 & 1.30 & 9.7 & 1.39\\
 & VSI & 6 & 0.4 & 1.09 & 0.6 & 4.60 & 3.5 & 1.21 & 4.7 & 1.30 & 6.5 & 1.43\\
\midrule
4K HEVC & frame & -- & 413.3 & 0.15 & 676.8 & 3.65 & 2,739.8 & 0.35 & 3,687.1 & 0.45 & 4,943.3 & 0.57\\
 & manifest & 1 & 22.3 & 3.35 & 31.1 & 6.86 & 99.9 & 3.53 & 131.5 & 3.64 & 173.3 & 3.75\\
 & manifest & 2 & 11.2 & 3.35 & 15.6 & 6.85 & 50.0 & 3.53 & 65.7 & 3.62 & 86.7 & 3.73\\
 & manifest & 4 & 5.6 & 3.33 & 7.8 & 6.84 & 25.0 & 3.49 & 32.9 & 3.62 & 43.3 & 3.74\\
 & VSI & 1 & 2.3 & 3.34 & 3.8 & 6.85 & 21.1 & 3.49 & 28.3 & 3.59 & 38.8 & 3.68\\
 & VSI & 2 & 1.1 & 3.33 & 1.9 & 6.84 & 10.6 & 3.47 & 14.1 & 3.57 & 19.4 & 3.67\\
 & VSI & 4 & 0.6 & 3.32 & 1.0 & 6.83 & 5.3 & 3.44 & 7.1 & 3.55 & 9.7 & 3.63\\
 & VSI & 6 & 0.4 & 3.34 & 0.6 & 6.86 & 3.5 & 3.46 & 4.7 & 3.58 & 6.5 & 3.65\\
\bottomrule
\end{tabular}
\end{table*}


Table~\ref{tab:stream} gives the result. Time is not the constraint. From the moment a segment closes, hashing its last fragment and signing take at most \NstreamLatMax\,ms, and the 99th percentile of the whole per-segment work, every fragment's hash and the provenance built and signed, is at most \NstreamPnnMax\,ms, against segments of one to six seconds. Counting all of that work, either C2PA method uses at most \NstreamCpuMax\% of one core; signing every frame uses up to \NstreamFrameCpuMax\%. These are feasibility figures from \NstreamSamplesMin{} or more segments per configuration, not a characterisation of tail latency under load. What the method decides is bandwidth. Under
ML-DSA-65 at 1080p, a manifest per one-second segment adds
\NstreamManOneKbps\,kbit/s, \NstreamManOneShare\% of the video; the session-key
method adds \NstreamVsiOneKbps\,kbit/s, \NstreamVsiOneShare\%, after a first
manifest of \NstreamVsiInit{} bytes; and signing every frame adds
\NstreamFrameOneKbps\,kbit/s, \NstreamFrameOneShare\% of the stream. Longer
segments divide the per-segment cost but delay authentication: a player can
check a segment only once it holds all of it, so verification lags capture by
at least the segment length. The session key here uses the signer's
algorithm, so the stream stays post-quantum end to end; a classical session
key under a post-quantum signer would shrink each segment's signature and
leave every segment forgeable by a quantum adversary.

A verifier that knows which stream it is playing replayed each stream with ten manipulations (Table~\ref{tab:streamdetect} in the supplement). Under both C2PA methods it detected a dropped, reordered, replayed or substituted segment, bytes truncated within a segment, and a segment taken from another stream by the same signer: truncation and substitution by the hash, the cross-stream segment by its stream identifier, and the rest by the sequence number. Without the initial manifest a session-key stream cannot be verified at all, and segments signed under an expired session key are rejected. One manipulation is not detected: removing the last segments. A stream that stops early looks like one whose tail was cut, unless its length is announced in signed data or the player gives up after a timeout; with an announced length, the cut is caught. A per-segment signature without a sequence number would also have let the drop, reorder and replay through. The earlier raw-frame experiment,
which hashed uncompressed YUV and signed every frame of 4K video for twenty minutes across ten configurations without a late frame or thermal throttling, remains in the supplement
(Table~\ref{tab:video}) as the computational upper bound.

\subsection{A verifier that reads its trust store from flash}
\label{sec:edge-verify}

A gateway or display device verifies against a local trust store. On the
Pi~5 we verified manifests from 500 signers with leaf and intermediate
embedded, and by thumbprint from a store held in memory and from the same
store after the page cache was dropped, so that every certificate came off
the SD card (Table~\ref{tab:edge} in the supplement). Once the certificates
are in memory, embedded and referenced verification cost the same to within
\NedgeWarmDiff\,ms, so carriage stays a question of bytes. A cold store is not
free on flash: reading the signer's certificates from the card adds
\NedgeColdEd\,ms to Ed25519 verification and \NedgeColdDSAsixtyfive\,ms to
ML-DSA-65.

\subsection{Credentials on a microcontroller}
\label{sec:mcu}

We signed and verified the same manifests on three microcontrollers and
measured the energy, and the radio cost of each carriage form over Bluetooth
LE. The Raspberry Pi Pico (RP2040, Cortex-M0+ at 125\,MHz, 264\,KiB of SRAM)
runs PQClean's portable C; the NUCLEO-L4R5ZI (STM32L4R5, Cortex-M4F,
640\,KiB) runs pqm4's assembly at 80\,MHz and at the 20\,MHz of pqm4's
benchmarks; the nRF52840-DK (Cortex-M4F at 64\,MHz, 256\,KiB) runs the same
objects and, for ES256 and SHA-256, its CryptoCell-310 through PSA Crypto.
Ed25519 and ES256 use Monocypher and micro-ecc. Each board holds one
manifest per algorithm in each carriage form; to verify, the firmware copies
it into a receive buffer as if it had just arrived, decodes it, rebuilds the
signed structure, checks the certificates up to the anchor in its trust store
and checks the manifest signature. Energy comes from a Nordic PPK2 sampling at
100\,kHz around marked batches of operations.

\begin{table}[t]
\caption{Signing and verifying on three microcontrollers, milliseconds, medians. \emph{Verify} is a complete manifest: receive it, decode it, check the leaf and intermediate up to the root anchor in the trust store, and check the manifest signature, with leaf and intermediate carried (\emph{L+I}) or only a thumbprint (\emph{Th.}). Pico: RP2040 Cortex-M0+ at 125\,MHz, portable C (PQClean). NUCLEO: STM32L4R5 Cortex-M4F at 80\,MHz, pqm4 assembly. nRF: nRF52840 Cortex-M4F at 64\,MHz, the same objects; \emph{ES256 (hw)} runs on its CryptoCell-310. Every board rejected a flipped signature bit, a tampered manifest, a tampered intermediate and a missing intermediate for every algorithm.}
\label{tab:mcu}
\small
\setlength{\tabcolsep}{4pt}
\begin{tabular}{lrrrrrrrr}
\toprule
 & \multicolumn{3}{c}{Pico} & \multicolumn{3}{c}{NUCLEO, 80\,MHz} & \multicolumn{2}{c}{nRF52840}\\
\cmidrule(lr){2-4}\cmidrule(lr){5-7}\cmidrule(lr){8-9}
Algorithm & Sign & L+I & Th. & Sign & L+I & Th. & Sign & L+I\\
\midrule
Ed25519 & 57.7 & 354.7 & 350.1 & 21.5 & 107.8 & 104.8 & 21.9 & 101.2\\
ES256 & 131.6 & 440.4 & 433.8 & 131.0 & 438.0 & 431.6 & 152.0 & 503.7\\
ES256 (hw) & -- & -- & -- & -- & -- & -- & 30.2 & 66.2\\
ML-DSA-44 & 154.8 & 160.7 & 143.3 & 96.3 & 130.9 & 112.1 & 90.5 & 121.4\\
ML-DSA-65 & 222.0 & 255.9 & 232.2 & 154.9 & 214.9 & 189.2 & 144.7 & 200.6\\
ML-DSA-87 & 329.7 & 412.5 & 381.5 & 190.0 & 360.1 & 325.7 & 199.6 & 333.1\\
\bottomrule
\end{tabular}
\end{table}


\paragraph{Time fits the budget}
On the slowest part, ML-DSA-44 and ML-DSA-65 verify faster than the classical schemes they replace, and ML-DSA-87 takes \NpicoChainDSAeightyseven\,ms (Table~\ref{tab:mcu}). The Pico checks an ML-DSA-44 manifest in
\NpicoChainDSAfortyfour\,ms and an ML-DSA-65 manifest in
\NpicoChainDSAsixtyfive, against \NpicoChainEd{} for Ed25519 and
\NpicoChainES{} for ES256. That comparison is between implementations as much
as algorithms, in the sense of Section~\ref{sec:backends}: Monocypher and
micro-ecc are compact portable libraries, and assembly-optimised Ed25519 for
the Cortex-M0 is faster than either~\cite{li2024m0}. At 20\,MHz the NUCLEO's
ML-DSA-65 signs in a median of \NnucTwentySignMcycles{} million cycles, in
line with pqm4's published counts~\cite{pqm4}. Signing takes
\NnucSignDSAsixtyfive\,ms on the NUCLEO at 80\,MHz and \NnrfSignDSAsixtyfive\,ms
on the nRF52840 for ML-DSA-65, a few frame periods at 30 frames per second,
and \NnucSignStackDSAsixtyfive\,KiB of stack.

Hardware support reverses the order. The CryptoCell-310 verifies an ES256
manifest in \NnrfChainEShw\,ms against \NnrfChainES{} in software on the same
chip, \NhwSpeedup{} times faster, which puts ES256 ahead of every ML-DSA
level: ML-DSA-44 takes \NhwRatioDSAfortyfour{} times as long and ML-DSA-87
\NhwRatioDSAeightyseven. The engine has no lattice support, so a product built
on this part gives up its accelerator when it gives up ECDSA.

\begin{table*}[t]
\caption{What a manifest costs a microcontroller beyond time. \emph{RAM}: received manifest, rebuilt signed structure and peak stack, all live at once, by carriage form (root, leaf and intermediate, leaf only, thumbprint); M0+ is the Pico's portable C, M4 the NUCLEO's pqm4 build. \emph{Verify}: energy for one complete verification, with leaf and intermediate carried, measured with the current vectors. The three boards are measured at different points (the whole Pico, the NUCLEO's MCU rail, the nRF52840's supply) and are not comparable with one another. Radio energy, measured, is in Table~\ref{tab:radio}.}
\label{tab:mcucost}
\small
\setlength{\tabcolsep}{3.6pt}
\begin{tabular}{lrrrrrrrr}
\toprule
 & \multicolumn{5}{c}{RAM (KiB)} & \multicolumn{3}{c}{Verify (mJ)}\\
\cmidrule(lr){2-6}\cmidrule(lr){7-9}
 & M0+ & \multicolumn{4}{c}{M4} & & & \\
\cmidrule(lr){3-6}
Algorithm & L+I & Root & L+I & Leaf & Th. & Pico & NUCLEO & nRF\\
\midrule
Ed25519 & 5.5 & 6.2 & 5.3 & 3.9 & 2.6 & 36.4 & 6.48 & 1.35\\
ES256 & 4.7 & 5.7 & 4.7 & 3.2 & 1.8 & 44.9 & 24.95 & 5.82\\
ES256 (hw) & -- & 9.9 & 8.9 & 7.4 & 6.0 & -- & -- & 0.89\\
ML-DSA-44 & 55.4 & 37.0 & 28.9 & 20.4 & 12.0 & 16.3 & 7.21 & 1.81\\
ML-DSA-65 & 83.3 & 47.8 & 36.8 & 25.3 & 13.8 & 25.9 & 11.84 & 3.01\\
ML-DSA-87 & 126.5 & 62.6 & 47.7 & 32.4 & 17.0 & 41.6 & 19.82 & 5.00\\
\bottomrule
\end{tabular}
\end{table*}


\paragraph{Memory and the radio}
Memory is where post-quantum credentials press on these parts
(Table~\ref{tab:mcucost} and Figure~\ref{fig:mcu}). On the NUCLEO, verifying
an ML-DSA-87 manifest needs \NramRootDSAeightyseven\,KiB with the root
carried, \NramChainDSAeightyseven\,KiB with leaf and intermediate,
\NramLeafDSAeightyseven\,KiB with the leaf alone and
\NramThumbDSAeightyseven\,KiB with a thumbprint; pqm4's verifier needs about
12\,KiB of stack, so the received and rebuilt buffers are the largest item.
The Pico's portable verifier needs \NramPicoDSAeightyseven\,KiB of its 264. On a
Cortex-M4 the carriage form, not the parameter set, decides whether a
verifier fits beside an application.

\begin{table}[t]
\caption{Measured Bluetooth LE energy for one manifest's COSE signature structure, with the certificates of each carriage form, between two nRF52840-DKs (2M PHY, 251-byte data length, 247-byte ATT MTU, 15\,ms connection interval, boards about a metre apart), on the measured board's supply, microjoules, medians of three rounds of ten transfers. \emph{Send}: from the first notification queued to the last acknowledged. \emph{Receive}: the same board as receiver, from the first notification received to the last. Both include the connected link's idle current and any link-layer retransmissions inside the window, not connection set-up or advertising. The receiving board checked every transfer byte for byte against its own copy; all arrived intact.}
\label{tab:radio}
\small
\setlength{\tabcolsep}{4pt}
\begin{tabular}{lrrrrrrrr}
\toprule
 & \multicolumn{4}{c}{Send} & \multicolumn{4}{c}{Receive}\\
\cmidrule(lr){2-5}\cmidrule(lr){6-9}
Algorithm & Root & L+I & Leaf & Th. & Root & L+I & Leaf & Th.\\
\midrule
Ed25519 & 205 & 164 & 87 & 14 & 214 & 169 & 89 & 14\\
ES256 & 228 & 173 & 91 & 13 & 242 & 180 & 94 & 14\\
ML-DSA-44 & 1,684 & 1,227 & 745 & 273 & 1,742 & 1,266 & 775 & 282\\
ML-DSA-65 & 2,310 & 1,669 & 1,017 & 366 & 2,387 & 1,713 & 1,045 & 380\\
ML-DSA-87 & 3,119 & 2,266 & 1,381 & 516 & 3,194 & 2,355 & 1,416 & 530\\
\bottomrule
\end{tabular}
\end{table}


We measured the radio both ways between two nRF52840-DKs, each manifest form
sent ten times per round in three rounds, with the measured board as sender
and then as receiver (Table~\ref{tab:radio}). Energy grows in a straight line
with the bytes: \NbleTxNj\,nJ a byte to send and \NbleRxNj\,nJ a byte to
receive (95\% intervals \NbleTxNjLo--\NbleTxNjHi{} and
\NbleRxNjLo--\NbleRxNjHi), fitted over \NbleTxN{} measured transfers from
\NbleTxXmin{} to \NbleTxXmax{} bytes with $R^2$ of \NbleTxRsq; the largest residual is \NbleTxResUj\,$\mu$J sending and \NbleRxResUj\,$\mu$J receiving, large only relative to the smallest, thumbprint-sized transfers. The windows cover the connection events
of a transfer, link-layer acknowledgements and any retransmissions, and the
connected link's idle current between them, but not connection set-up or
advertising. The radio pays for the carriage form directly. A receiver spends
\NrecvMChainDSAeightyseven\,mJ on an ML-DSA-87 manifest with leaf and
intermediate, \NrecvMLeafDSAeightyseven\,mJ on the leaf alone and
\NrecvMThumbDSAeightyseven\,mJ on a thumbprint, and receiving the chained
form costs \NrecvShareVerify\% of what verifying it costs on the same chip, so
a battery-powered verifier pays for the chain twice, once to receive it and
once to check it.

\begin{figure}[t]
\centering
\includegraphics[width=\linewidth]{figures/mcu}
\Description{Two grouped bar charts by algorithm. Left: RAM needed to verify
on the NUCLEO for four carriage forms; the ML-DSA bars fall from about 60 KiB
with the root carried to under 20 KiB with a thumbprint. Right: measured energy for an nRF52840 to receive
the manifest over Bluetooth LE in the same forms; ML-DSA-87 falls from about
3 mJ to about half a millijoule.}
\caption{Carriage form against verifier RAM on the NUCLEO-L4R5ZI
(Table~\ref{tab:mcucost}) and against the measured energy an nRF52840 spends
receiving the manifest (Table~\ref{tab:radio}).}
\label{fig:mcu}
\end{figure}

\subsection{Secure elements}
\label{sec:constrained}

The devices signing the most photographs are phones, which sign inside secure
elements, and no development board reproduces that. The measurements bound it
from both sides: an application processor signs every frame with time to
spare, and a Cortex-M4 needs a few frame periods and tens of kilobytes of
stack per ML-DSA signature. ML-DSA offers a way to split the work. The
external-$\mu$ variant lets the application processor hash the message
together with the public key into the 64-byte $\mu$ that ML-DSA signs, so a
secure element receives a fixed-size input whatever the message, and the
certificate profile allows it~\cite{fips204,rfc9881}. In a C2PA manifest the
signed message includes the protected header, and with it the embedded chain;
external $\mu$ keeps the chain's bytes out of the secure element, which then
does only the signing arithmetic instead of hashing tens of kilobytes streamed
to it over a slow bus. Whether a secure element keeps
up with burst capture remains to be measured.

%% ---------------------------------------------------------------------------
\section{Migration order and validation}
\label{sec:findings}

Cost is one half of a migration plan; the other is what the migrated credential
protects. This section asks which partial migrations leave which property of
Table~\ref{tab:properties} forgeable, and whether the reference implementation enforces the
rules a migration relies on.

\subsection{When the chain does not share one algorithm}
\label{sec:mixedchain}

Every chain so far uses one algorithm from leaf to root, and no migration
looks like that. A certificate carries two algorithms: the key it holds and
the signature its issuer put on it. Those come apart during a transition, and
in a particular order, because leaf certificates are reissued in days,
intermediates live for years, and roots are pinned in trust stores and move
last.

\begin{table}[t]
\caption{Chains whose levels migrate at different times, leaf first and root last; the root signs the intermediate but is not carried. \emph{Signer PQ} is whether every load-bearing signature in the signer's path is post-quantum. The last two columns add an RFC~3161 timestamp: with a classical TSA the time evidence stays forgeable, so a quantum adversary can backdate the manifest even when the signer's path is post-quantum.}
\label{tab:mixedchain}
\small
\setlength{\tabcolsep}{3.5pt}
\begin{tabular}{llrcrr}
\toprule
 & Chain & COSE & Signer & \multicolumn{2}{c}{+ timestamp (B)}\\
\cmidrule(lr){5-6}
Stage & leaf/int./root & (B) & PQ & ES256 TSA & ML-DSA-65 TSA\\
\midrule
Classical & ES256/ES256/ES256 & 1,585 & no & 3,651 & 17,209\\
Leaf migrated & D65/ES256/ES256 & 6,724 & no & 8,788 & 22,348\\
Leaf and intermediate & D65/D65/ES256 & 11,867 & no & 13,933 & 27,491\\
PKI first, leaves last & ES256/D65/D65 & 9,971 & no & -- & --\\
Fully migrated & D65/D65/D65 & 15,113 & yes & 17,179 & 30,737\\
Assurance ladder & D44/D65/D87 & 14,902 & yes & -- & --\\
Ladder inverted & D87/D65/D44 & 16,182 & yes & -- & --\\
Cheap leaves, strong root & D44/D44/D87 & 13,373 & yes & -- & --\\
\bottomrule
\end{tabular}
\end{table}


Table~\ref{tab:mixedchain} separates two things usually spoken of as one.
Bytes are set by which levels have migrated. Quantum resistance is set by the
weakest load-bearing signature: the manifest signature, the signature on the
leaf certificate and the signature on the intermediate. The root's
self-signature is not load-bearing, and the root is not carried, but its
algorithm still matters because it signs the intermediate. A post-quantum
leaf under the existing classical PKI costs \NmixLeafMig{} times the
classical structure, \NmixPctFullLeafMig\% of what full migration costs;
migrating the intermediate too brings it to \NmixLeafInt{} times. Neither
makes the signer's path quantum resistant, because an adversary who can forge
the issuer's ECDSA signature mints a leaf certificate and signs whatever it
likes.

Migrating the PKI first while issuing classical leaves for compatibility is
the order an operator reaches by reasoning about disruption rather than about
adversaries. It costs \NmixPkiFirst{} times the classical structure,
\NmixPctFullPkiFirst\% of the full price, and leaves the manifest itself
forgeable, since its signature is still ECDSA. What it does buy is a
post-quantum trust infrastructure ready to issue post-quantum leaves without a
second change at the anchor, so its value lies in what follows it, not in any
protection it gives the manifests signed meanwhile.

The timestamp needs its own migration. With a classical time-stamping
authority, a fully migrated signer's manifest can still be backdated by an
adversary who forges the timestamp, which undoes the argument a timestamp
exists to make, that the manifest was signed before its certificate was
revoked. A post-quantum timestamp costs \NmixWholeFullPq{} bytes of manifest
cryptography against \NmixWholeFullEs{} with an ES256 authority, and the
provenance is quantum resistant only when both are post-quantum.

An assurance ladder, a category-2 leaf under a category-3 intermediate under a
category-5 root, puts the higher margin on the longer-lived key and costs
\NmixLadder{} times the classical structure, against \NmixFull{} for a uniform
category-3 chain: the strongest key sits in the root, which is not carried.
Inverting the ladder costs \NmixInverted{} times and is the configuration to
check for in review, since it puts the weakest parameter set on the key that
lives longest and is hardest to replace.


\subsection{What the reference implementation accepts}
\label{sec:validation}

\begin{table*}[t]
\caption{Validation through the reference implementation (\texttt{c2pa-rs} \texttt{e2ad88e} with our ML-DSA patches, p2): a selection of the 40 cases; Table~\ref{tab:valfull} in the supplement lists all. T, V, I: Trusted, Valid, Invalid; E: reader error; SR: signing refused (\emph{direct}: our harness's own structure with that certificate). \emph{Expected} follows C2PA~2.5. \emph{TST}: the timestamp's status (t trusted, m mismatch, u untrusted); an ignored timestamp leaves the signer judged now. The TSA is ES256 except in the indented rows. $^\dagger$Bold: a defect, also present unpatched. The CA-signed \emph{revoked} OCSP row is where C2PA~2.5 and RFC~6960 disagree (Section~\ref{sec:validation}).}
\label{tab:validation}
\footnotesize
\setlength{\tabcolsep}{3pt}
\begin{tabular}{llccc}
\toprule
 & & ES256, & ML-DSA & \\
Case & Expected & Ed25519 & 44/65/87 & TST\\
\midrule
Leaf and intermediate & T & T & T & --\\
\quad timestamped & T & T & T & t\\
Leaf, intermediate, root & T & T & T & t\\
Leaf only, issuer is anchor & T & T & T & t\\
Leaf only, intermediate missing & not T & V & V & t\\
Unknown anchor & not T & V & V & t\\
Thumbprint (\texttt{x5t}) & not T & E & E & --\\
Leaf EKU serverAuth & refused & SR (direct I) & SR (direct I) & -- (direct --)\\
Leaf EKU timeStamping only$^\dagger$ & not T & \textbf{T} & \textbf{T} & t\\
COSE algorithm $\neq$ key (ML-DSA-44 only) & I & I & I & --\\
Image changed after signing & I & I & I & t\\
Signature bit flipped & I & I & I & m\\
Expired leaf, timestamp inside validity & T & T & T & t\\
Expired leaf, no timestamp & not T & I & I & --\\
Expired leaf, timestamp over other data & not T & I & I & m\\
Expired leaf, TSA not trusted & not T & I & I & u\\
Timestamp over other data & T & T & T & m\\
TSA not trusted & T & T & T & u\\
TSA certificate lacks timeStamping$^\dagger$ & T & \textbf{I} & \textbf{I} & u\\
ML-DSA-65 TSA & T & T & T & t\\
\quad timestamp over other data & T & T & T & m\\
\quad expired leaf, timestamp inside validity & T & T & T & t\\
OCSP, CA-signed, revoked & T & T & T & t\\
OCSP, delegated responder, revoked & not T & I & I & t\\
\bottomrule
\end{tabular}
\end{table*}


Table~\ref{tab:validation} gives a selection of the \Nvalcases{} cases of
Section~\ref{sec:patch}; Table~\ref{tab:valfull} in the supplement lists every
case for every algorithm. With both patches every ML-DSA case behaves as the
classical ones do, and with the signing patch alone none reaches
\emph{Trusted}, because path validation cannot check an ML-DSA certificate.

\paragraph{A fully post-quantum credential works}
A credential with an ES256 time-stamping authority validates in
\NvalVerifyLo{} to \NvalVerifyHi\,ms on the M1 Pro. An ML-DSA-65 signer,
timestamped by an ML-DSA-65 time-stamping authority under an ML-DSA-65 anchor,
validates as \emph{Trusted} with a trusted timestamp in \NvalPqTsaMs\,ms, and still does after the signer's certificate
has expired, which is the job a timestamp exists to do. A timestamp over other
data, or from an untrusted authority, is ignored as the specification
says~\cite{c2pa25}, and an expired signer with such a timestamp is rejected, so
a rejected timestamp never lends a signer historical validity. The library
rejects a modified image, a flipped signature bit and a COSE algorithm that
does not match the key, stops short of \emph{Trusted} for a chain to an unknown
anchor or a leaf whose intermediate is missing, refuses to sign with a leaf
that fails its own certificate profile, and accepts the leaf alone when the
leaf's issuer is an anchor, the configuration Section~\ref{sec:carrying}
recommends. It does the same for the time-stamping authority: a token that
carries only the authority's certificate yields a trusted timestamp when the
authority's issuing CA is a time-stamping anchor, for every algorithm, and an
untrusted one when only the root is. A manifest carrying a thumbprint fails with a reader error rather
than a validation status, a 50-certificate chain is reported \emph{Valid} but not \emph{Trusted}, and a
1\,MB certificate validates in \NvalBigCertMs\,ms using \NvalBigCertMB\,MB of
memory.

\paragraph{A time-stamping certificate is accepted as a claim signer}
A leaf whose only extended key usage is \texttt{id-kp-timeStamping} validates
as a \emph{Trusted} claim signer, for every algorithm, and still does under a
policy whose usage list names only \texttt{c2pa-kp-claimSigning}. C2PA~2.5
says that a certificate valid for time-stamping shall be valid for that purpose
alone~\cite{c2pa25}. The library's key-usage check accepts e-mail protection,
time-stamping and OCSP signing unconditionally, whatever the configured list
says. Two added cases confirm it for every algorithm, through the tagged
release for ES256 and Ed25519 and through the patched library for ML-DSA: a
leaf valid only for e-mail protection, or only for OCSP signing, also
validates as a \emph{Trusted} claim signer under the claim-signing-only
policy, and the OCSP-signing leaf does so even though C2PA~2.5 confines such a
certificate to that one purpose. Whoever holds the key of a time-stamping
authority whose certificate chains to a claim-signing anchor can therefore
sign manifests that a validator using the library's key-usage check reports as
\emph{Trusted}. Such a manifest is attributed to the authority's own subject
name, not to an arbitrary publisher: the defect lets a key operated for
another purpose, often by another party, produce trusted provenance, which
the separation the specification draws is meant to prevent, but it does not by
itself let that key impersonate another signer. On the live lists the
condition is not hypothetical, and we checked it with real certificates. We
pulled the time-stamping certificates out of the RFC 3161 tokens of the
production files (Section~\ref{sec:prodchains}), \NtpcChains{} distinct chains
in all, and ran each through the library's own certificate-profile and trust
checks against the C2PA claim-signing list, with no signature. Both checks pass
for \NtpcStrong{} of the \NtpcChains{} chains as a claim signer under a
claim-signing-only policy: \NtpcGoogle{} of Google's time-stamping leaves and
vivo's, the ones whose issuing CA is on the C2PA time-stamping list and is
itself issued by a claim-signing anchor. These \NtpcStrong{} chains sign the
timestamps in \NtpcFiles{} of the production files. For each, the profile
check accepts the time-stamping-only EKU and the trust check builds a path to a
claim-signing anchor, so the library would report a manifest signed by the
key behind any of them as \emph{Trusted}; only the private key is absent. The
remaining chains fail because their issuer is not on the claim-signing list
(DigiCert's, Sectigo's and other general-purpose time-stamping authorities) or,
for \NtpcCaRej{} of them, because the profile check correctly refuses a CA
certificate. The two trust lists make the overlap structural rather than
incidental: \NtlBothLists{} of the \NtlClaimRoots{} roots on the claim-signing
list are also roots on the time-stamping list, and all \NtsaIcaUnderClaim{}
time-stamping issuing CAs on the time-stamping list are issued by claim-signing
roots, so any time-stamping authority listed under one of them is in the same
position as Google's. The behaviour reproduces in the unmodified library at commit
\texttt{\NvalCommit}, in the latest tagged release, c2pa-v0.91.9, and in the
command-line tool c2patool~0.28.4, so it is not an artefact of our patch;
Table~\ref{tab:evidence} in the supplement records the production-certificate
check and what else was reproduced where.

\paragraph{A timestamp the validator should ignore invalidates the manifest}
In the other direction, a time-stamping authority whose certificate lacks the
time-stamping usage makes the whole manifest \emph{Invalid}, with status
\texttt{signingCredential.invalid}, where the specification says to report the
timestamp as untrusted and judge the signer at the current time. The token
travels in the COSE unprotected header, so anyone who can rewrite a file can
swap in such a token, and a misconfigured authority produces one by accident.
Either way valid content is reported as \emph{Invalid}, the status of
tampered content, instead of falling back to validation at the current time.

\paragraph{Revocation stapled in a CA-signed response has no effect}
Unlike the two defects above, this behaviour arises where the C2PA profile and
RFC~6960 disagree and the implementation follows the profile.
RFC~6960 lets the issuing CA sign an OCSP response itself, with no special key
usage~\cite{rfc6960}. C2PA~2.5 instead requires a response to be signed by an
end-entity certificate valid for \texttt{id-kp-OCSPSigning}, while also asking
validators to check stapled responses as RFC~6960 does, which accepts the
CA-signed form~\cite{c2pa25}. The reference implementation resolves the
tension by ignoring the response: it looks for the responder's certificate in
the response, finds none, and treats the status as unknown, the same outcome as
when nothing is stapled. A signer whose stapled, CA-signed response says
\emph{revoked} therefore validates as \emph{Trusted}
(Table~\ref{tab:validation}), with online OCSP fetching off as it was for the
whole matrix; a validator that also fetches status could still catch it. The risk is to
deployments that believe their stapled evidence protects them. A generator that
staples revocation status, as \NcorpProdOcsp{} production signatures do, must
use a delegated responder for it to count; under ML-DSA-65 a delegated
response costs \NocspDelDSAsixtyfive{} bytes per manifest, because it carries
the responder's certificate, against \NocspCADSAsixtyfive{} for a CA-signed
one. Golaszewski et al.\ found that
revocation checking is optional in C2PA~\cite{golaszewski2026}; here evidence
is present and ignored.

\paragraph{Disclosure}
I reported the invalidating-timestamp behaviour publicly, with minimal
vectors, as c2pa-rs issue \#3005 on 3 November 2026~\cite{c2pars3005}. The
key-usage behaviour lets a time-stamping key sign trusted manifests, so I
reported it privately, with its vectors, to Adobe's product security team,
which handles security reports for the library, on 9 November 2026. Its cause
and its reach are described here, but its signed vectors are withheld from the
released data until a fix ships or 90 days after the report, whichever comes
first; its status codes are released. The OCSP behaviour follows the
specification, so I raised it with the specification's
editors~\cite{c2paspec143}, as I did the reading behind leaf-only
carriage~\cite{c2paspec142}.

\paragraph{Reading the findings}
The findings are two implementation defects, the key-usage and
invalidating-timestamp behaviours, and one interoperability and
revocation-policy issue, the CA-signed OCSP response that the C2PA profile and
RFC~6960 treat differently. None is specific to post-quantum algorithms; each
appears with ES256 and Ed25519. Each, however, breaks a property a migration
relies on: that a time-stamping key cannot attribute content, that a timestamp
is judged as the specification says, and that a revoked signer stays revoked.
Several of the Content Authenticity Initiative's SDKs for other languages reuse
this library, so the findings may affect those integrations too, though we did
not evaluate any application built on it or measure deployment share. The two
defects are of the kind that adversarial testing of TLS libraries exposed a
decade ago~\cite{brubaker2014frankencerts,debnath2021ceres}. They surfaced
because adding an algorithm called for a conformance matrix with negative
cases, and a migration that adds the algorithm without one carries such defects
over unnoticed.

%% ---------------------------------------------------------------------------
\section{Discussion and limitations}
\label{sec:discussion}

\paragraph{Where the standards stand}
The measurements were planned against C2PA~2.4, which permitted only
classical algorithms, and C2PA~2.5 added ML-DSA while they were under way.
The specification has therefore caught up with RFC~9964 and the reference
implementation has not: at commit \texttt{\NvalCommit} of 2~November 2026,
\texttt{c2pa-rs} signs only classical algorithms, and making it validate an
ML-DSA chain as trusted took the four-gate patch of Section~\ref{sec:patch}
and a second one for path validation. Every size in this paper is a property of the encoding and carries over to any conforming implementation up to a constant of its own, which for \texttt{c2pa-rs} is about \Nlibgap{} bytes, mostly padding (Section~\ref{sec:patch}); every
timing is a property of the libraries named beside it. Nothing in C2PA~2.5
changes the carriage rules this paper relies on, and its trust-list snapshot
identifiers are what a cache needs to invalidate correctly
(Section~\ref{sec:cache}).

\paragraph{What the recommendation asks of whom} Generators can stop carrying the root today. Leaf-only carriage needs the signer's issuing CA on the trust list, which is a decision for CAs and for the trust-list operator, and it shortens the path a validator checks as well as the manifest. It needs no change to the specification, but it is a change of trust policy, and the operator takes on four things in exchange for the saving. An anchor is trusted because it is listed: a validator that treats an issuing CA as an anchor no longer checks it against its root's revocation status, so withdrawing that CA means republishing the list, and every verifier must pick the change up. Issuing CAs rotate far more often than roots, so the list changes more often, and offline verifiers, which hold a copy, see each change later. Constraints that a root imposes on the CAs below it, on names, policies or path length, stop applying unless the anchor entry carries them. And an issuing CA retired from the list leaves its older manifests verifiable only by a validator that kept it, or that falls back to the full chain, which a generator should therefore still be able to emit. The trust-list snapshot identifiers of C2PA~2.5 make each change addressable, and Section~\ref{sec:cache} shows why a validator must honour them. Validators should tie cached trust decisions to the trust list, policy and validation time they were made under, and bound them by certificate validity and the currency of revocation status. The specification could say both outright.

\paragraph{Limitations}
Remote trust-list resolution is not measured; it is a network round trip
whose cost belongs to the deployment. The post-quantum certificates are ours,
issued under a production-like profile, since no deployed signer issues them
yet; production profiles vary by a few hundred bytes per certificate
(Section~\ref{sec:method-decisions}), which moves the ratios we report by up to about ten per cent. Falcon and the composites carry placeholder object
identifiers, and Falcon is measured as submitted rather than as FN-DSA will
be standardised. HQC is not measured, and neither is any key establishment;
encrypted archives, whose data keys would have to be re-wrapped under ML-KEM,
lie outside the paper's question, and Section~S1 of the supplement measures
that re-wrap separately. The validation matrix covers an ML-DSA-65 time-stamping authority only at that one parameter set, and every timestamp in it comes from our own authority, not a deployed one. The microcontroller energy figures are taken at three
different points and compare algorithms on one board rather than boards with
one another, and the radio figures are for one Bluetooth LE configuration, two boards about a metre apart in an office, with whatever interference that brought. No secure element and no accelerator beyond the CryptoCell was
measured, and batched GPU verification of lattice
signatures~\cite{shen2024gpu,lee2024falcon} is a question we leave open.
Timings on the Xeon come from a small, two-core cloud virtual machine whose host load we do not control; the across-run intervals record what that cost, and its scaling curve stops at two workers.

%% ---------------------------------------------------------------------------
\section{Related work}
\label{sec:related}

\subsection{Certificate size in post-quantum TLS}
\label{sec:related-tls}
That certificates rather than signatures dominate a post-quantum exchange is
established in the TLS literature, and this paper does not claim it as new.
Sikeridis et al.\ found the certificate chain, not the handshake signature,
driving the added latency of post-quantum TLS~1.3, to the point of proposing
different algorithms at different levels of one
chain~\cite{sikeridis2020ndss,sikeridis2020conext}; later studies reached the
same conclusion on newer stacks~\cite{sosnowski2023}, and Cloudflare made the
sizing argument to practitioners~\cite{westerbaan2021}. The remedies all stop
sending certificates the verifier could already have: cached
information~\cite{rfc7924}, abridged certificate compression, which ships a
dictionary of the WebPKI's roots and intermediates so that only the leaf
travels and reports about 75\% off the chain against about 25\% for generic
compression~\cite{certabridge}, and Merkle Tree
Certificates~\cite{merkletreecerts}. Leaf-only carriage is the manifest
equivalent of abridged compression in its simplest form, with the C2PA trust
list as the dictionary, and we present it as such.

\paragraph{What transfers and what does not}
The structural result transfers, and our measurements extend it to Falcon,
SLH-DSA and composites: because every certificate contains a signature, chain
dominance is a property of chains rather than of any scheme. Generic
compression fails here as it does in TLS (Table~\ref{tab:compress}). Four
differences change the conclusions. A manifest chain is paid once per object
and kept for the life of the file, so the cost is storage and bandwidth rather
than handshake latency. Nothing is negotiated, so the carriage form is a
policy for every verifier. A manifest verifier may be offline for good, which
is why we recommend the leaf over the thumbprint where the TLS analogy would
favour the smaller option. And C2PA's trust list lists issuing CAs as anchors,
which makes the leaf-only form conformant without any new mechanism.

\subsection{Post-quantum benchmarking}
\label{sec:related-bench}

Primitive-level performance is well covered, from pqm4 on
microcontrollers~\cite{kannwischer2019pqm4,pqm4,greconici2021,li2024m0} and
NEON and AVX2 implementations on application
processors~\cite{becker2022neon} to batched GPU verification reported in the
millions per second~\cite{shen2024gpu,lee2024falcon,nvidia2024cupqc}. At the
protocol level, Paquin et al.\ benchmarked post-quantum TLS under emulated
loss~\cite{paquin2020}, and the Open Quantum Safe project made such
experiments routine~\cite{stebila2016}. Deployment has followed for key
agreement and not for authentication: by late 2025 most human-initiated
traffic to Cloudflare used post-quantum key agreement while no publicly
trusted post-quantum certificate was in use~\cite{westerbaan2025pq}. The half
of the migration this paper prices has not begun even on the web.

Methodology has become a subject of its own. Riou et al.\ catalogue pitfalls
in ML-DSA benchmarking, among them the tail that rejection sampling gives
signing time~\cite{riou2026}; SUPERCOP has long recorded the compiler and
options with every measurement~\cite{ebacs}, and Mytkowicz et al.\ showed that
measurement setup alone can reverse a conclusion about \texttt{-O2} against
\texttt{-O3}~\cite{mytkowicz2009}. Our audit (Section~\ref{sec:audit}) finds
that applied post-quantum work rarely carries those habits over.
Implementation-controlled measurement is rarer still; our earlier study of HQC
as a fallback to ML-KEM carries such a control for key
establishment~\cite{alfarajat2026hqc}. It shares this paper's method and not
its data or subject: its hardware work runs key establishment over a Thread
link, and the microcontroller work here is signatures only, over Bluetooth LE,
on separate firmware.

\subsection{Content provenance}
\label{sec:related-prov}

C2PA manifests are JUMBF boxes~\cite{iso19566-5} carrying COSE signatures, and
they have left the laboratory: Leica shipped the first camera that signs at
capture in 2023~\cite{cai2023leica}, Sony added signing to its bodies the
following year~\cite{sony2024c2pa}, and Samsung and Google brought content
credentials to phones in 2025~\cite{samsung2025s25,google2025pixel}. Google's
design keeps signing keys in a dedicated security chip, runs a time-stamping
authority on the device and issues a certificate for every
image~\cite{google2025pixel}, the lifecycle under which Section~\ref{sec:cache}
finds caching worth least.

Golaszewski et al.\ give the first independent security analysis
of C2PA and conclude that the specification does not meet its stated
goals~\cite{golaszewski2026eprint,golaszewski2026}. Two of their findings
concern machinery whose post-quantum cost we measure: revocation checking is
optional, and generators and validators disagree about which timestamp is
trusted. The two departures we found in the reference implementation
(Section~\ref{sec:patch}) are of the same kind, in the key-usage rules for
signers and time-stamping authorities. VerITAS proves in zero knowledge that a
published photograph differs from a signed original only by permitted
edits~\cite{datta2025veritas}, and Nemecek et al.\ build images whose valid
manifests contradict their watermarks~\cite{nemecek2026integrity}. Watermarks
and perceptual hashes are the alternatives to a signature; they rest on no
assumption Shor's algorithm breaks, and they fail by removal or
steering~\cite{wen2023treering,gowal2025synthid,an2024waves,zhao2024invisible,%
meta_pdq,photodna,struppek2022,jain2022}.

Post-quantum signatures have been proposed for image authenticity outside this
ecosystem, with zero-knowledge proofs of permitted
transformations~\cite{iavich2025pqveritas} and inside a watermarking framework~\cite{alkhatib2025deepfake}. Neither measures the credential, and to our knowledge this paper is the first published measurement
of post-quantum signatures inside the credential stack that cameras and phones
ship.

\subsection{Migration measurement}
\label{sec:related-migration}

Mixed chains have been measured in TLS: Paul et al.\ placed different
signature algorithms at different levels of a TLS~1.3 chain~\cite{paul2022},
and Bindel et al.\ tested whether hybrid certificates survive real X.509, TLS
and S/MIME implementations~\cite{bindel2017}. Section~\ref{sec:mixedchain}
asks what each ordering protects and what it costs to carry in every object.
Planning for migration is mostly guidance: research agendas~\cite{ott2019},
crypto-agility~\cite{nist_cswp39}, inventory~\cite{nist_sp1800_38,ecma424} and
maturity models~\cite{hohm2023}. Past migrations were slow at the tail: four
years after the MD5 rogue CA~\cite{stevens2009} a trusted CA still issued
MD5-signed certificates~\cite{durumeric2013}, and TLS clients abandon old
algorithms slowly~\cite{kotzias2018}. Provenance has a worse tail than TLS: a
handshake renegotiates on the next connection, and a signed image keeps its
signature until someone signs it again.

%% ---------------------------------------------------------------------------
\section{Conclusion}
\label{sec:conclusion}

Making media provenance quantum-safe costs less than the size of a
post-quantum signature suggests and more than a signature swap: a manifest
under ML-DSA-65 is \NvsESDSAsixtyfive{} times a matched ES256 manifest and \NvsPSDSAsixtyfive{} times a matched PS256 one, the two algorithms deployed signers use, and a timestamped one
\NtsVsESDSAsixtyfive{} and \NtsVsPSDSAsixtyfive{} times, because certificates
were already most of the structure and every certificate carries a
signature of its own. The finding holds from Falcon to SLH-DSA and for
composite signatures, and generic compression does not touch it.

How much of the chain travels is therefore the design question. Leaving the
root out, as the specification says generators should and most signers already do,
saves a quarter of the manifest for those that still carry it. Carrying the leaf alone is conformant when the
leaf's issuer is on the trust list, saves a further third, and halves the
path a validator checks; it needs CAs to list their issuing CAs as anchors,
which several already do. The thumbprint is smaller still and neither
conformant nor verifiable offline.

Time is not the obstacle on the application processors and workloads we measured. On a phone photograph, the hash the format already requires costs
more than the post-quantum signature; on an edge board, signing an encoded 4K
stream by either C2PA live method uses a fraction of one core, and the method,
not the algorithm, sets the bandwidth. On a microcontroller the verifier's
memory and the radio are the constraints, and the carriage form decides both.

Two findings concern how the field measures. An apparent architecture effect
in our first results came from a Python package whose macOS OpenSSL was built
without optimisation, now fixed, and only \NauditBothWord{} of \NauditNWord{} applied studies we audited state both the compiler and the optimisation level their numbers depend on;
and a capacity estimate built on median signing times is optimistic by about
a fifth for ML-DSA, because rejection sampling gives signing a tail that a
batch pays in full. A post-quantum timing needs the implementation, its
version, its build and a measure of spread.


\section*{Data and code availability}
The measurement code, the firmware, every result file, the scripts that turn
them into this paper's tables and numbers, the validation harness with its
patches and test vectors, and the corpus inventory are released at
\url{https://github.com/moathfarajat-commits/chain-not-signature} and archived as
\url{https://doi.org/10.5281/zenodo.24091529} (release 1.0.0). The production-signed files
themselves are not redistributed, since several of their sources do not
allow it; the inventory gives each file's source URL, licence and SHA-256.
The signed vectors for the key-usage behaviour are withheld for the period
Section~\ref{sec:validation} gives.

\begin{acks}
The author received no funding for this work. Claude, an AI model by
Anthropic, was used to generate the LaTeX code of this article and to help
with its writing; the author reviewed and edited all of the text and takes
full responsibility for its content.
\end{acks}

\bibliographystyle{ACM-Reference-Format}
\bibliography{refs}

\end{document}
